\documentclass[10pt]{amsart}
\usepackage[a4paper,margin=23mm]{geometry}
\usepackage{amsmath,amssymb,amsthm,mathtools,hyperref,graphicx,xspace,letltxmacro,xpatch}
\usepackage[english]{babel}
\usepackage{xfrac}

\usepackage[scr=rsfs,cal=euler]{mathalfa}
\multlinegap0pt
\hyphenation{another exists inequal-i-ties inequal-i-ty inter-pre-ted obtained}

\newenvironment{enumeratei}
{\bgroup\def\theenumi{\roman{enumi}}\def\theenumii{\arabic{enumii}}\begin{enumerate}}
{\end{enumerate}\egroup}
\newenvironment{enumerateA}
{\bgroup\def\theenumi{\Alph{enumi}}\def\theenumii{\arabic{enumii}}\begin{enumerate}}
{\end{enumerate}\egroup}


\newcommand{\Psfrac}[2]{(\sfrac{#1}{#2})}
\newcommand{\psfrac}[2]{\sfrac{(#1)}{#2}}
\newcommand{\spfrac}[2]{\sfrac{#1}{(#2)}}
\newcommand{\pspfrac}[2]{\sfrac{(#1)}{(#2)}}
\newcommand{\Ppsfrac}[2]{(\sfrac{(#1)}{#2})}

\newcommand{\rcount}{\mathrm{count}}
\newcommand{\norm}{\mathrm{norm}}
\newcommand{\rG}{\mathrm{G}}
\newcommand{\rS}{\mathrm{S}}
\newcommand{\rV}{\mathrm{V}}
\newcommand{\rL}{\mathrm{L}}
\newcommand{\rO}{\mathrm{O}}

\theoremstyle{plain}
\newtheorem{theorem}{Theorem}[section]
\newtheorem{proposition}[theorem]{Proposition}
\newtheorem{corollary}[theorem]{Corollary}
\newtheorem{lemma}[theorem]{Lemma}

\theoremstyle{definition}
\newtheorem{definition}[theorem]{Definition}
\newtheorem*{claim}{Claim}
\newtheorem{assumption}[theorem]{Assumption}
\newtheorem{example}{Example}[section]
\newtheorem*{remark}{Remark}
\newtheorem*{notation}{Notation}

\let\eps\varepsilon

\newcommand{\RR}{\mathbb{R}}
\newcommand{\NN}{\mathbb{N}}
\newcommand{\ZZ}{\mathbb{Z}}
\newcommand{\cC}{\mathcal{C}}
\newcommand{\drm}{\mathrm{d}}
\newcommand{\euler}{\mathrm{e}}
\newcommand{\cD}{\mathcal{D}}
\newcommand{\cE}{\mathcal{E}}
\newcommand{\cM}{\mathcal{M}}
\newcommand{\cA}{\mathcal{A}}
\newcommand{\cQ}{\mathcal{Q}}
\newcommand{\cF}{\mathcal{F}}
\newcommand{\cG}{\mathcal{G}}
\DeclareMathOperator{\supp}{supp}
\DeclareMathOperator{\Deg}{Deg}
\DeclareMathOperator{\spec}{spec}
\DeclareMathOperator{\arsinh}{arsinh}
\DeclareMathOperator{\Eins}{\mathbf{1}}

\let\oldint\int
\renewcommand{\int}{\oldint\limits}

\usepackage{xfrac,xcolor,graphicx}\graphicspath{{./figures/}}\providecommand{\skpt}{\par\smallskip\noindent}\providecommand{\pointir}{.\ }\providecommand{\eg}{e.g.}\providecommand{\ie}{i.e.}\providecommand{\resp}{resp.}\providecommand{\cf}{cf.}\providecommand{\etc}{etc.}\let\Subsection\subsection\let\Subsubsection\subsubsection
\title{Gaussian heat-kernel upper bounds characterized by Sobolev inequalities and volume doubling on irregular graphs}
\author{Matthias Keller, Christian Rose}
\date{doi:10.5802/jep.340; CC BY 4.0}
\begin{document}
\maketitle
\noindent\textit{Editorial note.} Complete original author mathematical body, all local proofs and appendices retained; all original assumptions/ranges remain unchanged. Publisher bibliography restored to native citation keys, including same-year entries matched to original PDF citation order. Only wrapper and cedram matter directives/whitespace adapted.\par\medskip
\noindent\textit{Allowed standard prerequisites (editorial).}\begin{itemize}
\item Standard functional analysis, measure theory and the author-cited classical theorems. Standard prior results stay cited; complete source-local proofs and definitions retained.
\end{itemize}\medskip
\vspace*{-1.5\baselineskip}
\tableofcontents

\section{Introduction and main results}

Geometric characterizations of upper heat kernel bounds go back to the work of Varopoulos \cite{Varopoulos-85}. There it was shown that uniform diagonal upper heat kernel bounds of Dirichlet forms are characterized by uniform Sobolev inequalities. Different characterizations in terms of Nash, Gagliardo-Nirenberg, and Faber-Krahn inequalities have been obtained by several authors over the last decades in different settings, \cite{CarlenKS-87,Coulhon-92,SaloffCoste-92,Grigoryan-94,
SaloffCoste-92a,SaloffC-01,BCS, GrigoryanHH}. In~the case of continuous-time heat kernels on graphs, strong bounds on the underlying geometry have been studied since the seminal work of Davies and Pang \cite{Davies-93,Pang-93}.
Among many other milestones, Delmotte's fundamental work on the characterization of full Gaussian bounds for heat kernels with respect to the normalizing measure is mentioned \cite{Delmotte-99}. Furthermore, the textbook \cite{Barlow-book} provides characterizations of heat kernel upper bounds on graphs with strong boundedness assumptions on the underlying geometry.

In the case of graphs with unbounded geometry only partial results exist. The Davies-Gaffney-Grigor'yan estimate, an integrated heat kernel bound, has been obtained on graphs in \cite{BauerHuaYau-17}. Delmotte's work has been extended to graphs with normalizing measure but possibly unbounded combinatorial geometry in \cite{BarlowChen-16}. Graphs with bounded combinatorial geometry but possibly unbounded weights are studied in \cite{AndresDS-16,Bella-22} which is related to work in the continuum of Trudinger \cite{Trudinger-71}. Grigor'yan's two-point method to obtain off-diagonal from on-diagonal bounds has been extended to graphs with unbounded geometry involving intrinsic metrics \cite{Folz-11}. In~\cite{KellerRose-22b,KellerRose-22a} the authors obtained Gaussian upper bounds for large times assuming Sobolev and volume doubling properties on large scales in terms of intrinsic metrics.

Here we are interested in the geometric characterization of Gaussian upper bounds for the heat kernel on graphs with unbounded geometry.
We provide an equivalence of scale-invariant
Sobolev inequalities and the conjunction of Gaussian bounds, volume doubling property on large scales, and a local regularity property.
To allow for more unboundedness in the geometry than in the classical setting we have to expand on well known concepts using various ideas.

The first one is that we allow for a variable dimension in the Sobolev inequality. While this dimension function can possibly be unbounded in general, we~can relate it to the volume doubling dimension and a growth rate in the case where the vertex degree is polynomially bounded.

A second idea concerns the derivation of a volume doubling property from the Sobolev inequality. Due to the discrete structure of the space additional error terms are unavoidable. However, they can be controlled for large radii.

The third idea is a local version of regularity property appearing already in a uniform form in \cite{BarlowChen-16} for the normalizing measure. Indeed, this property is an immediate consequence of the Sobolev inequality by plugging in characteristic functions. However, it allows to formulate all occurring bounds in terms of the vertex degree at centers of balls. This way this local regularity property becomes part of the characterization.

Finally, we~extend the Gaussian bounds developed in \cite{KellerRose-22a} derived from volume doubling and Sobolev inequalities. The correction terms now depend only on degrees of the vertices at which the heat kernel is evaluated. We further emphasize that the off-diagonal bounds are indeed much finer than in previous works and are sharp in specific situations such as \cite{Davies-93,Pang-93}.

Although the main focus of the work is graphs of unbounded geometry, already in the special case of the normalized Laplacian, Theorem~\ref{thm:equivnormalized}, our result is new as it requires much less boundedness assumptions on the geometry than earlier works.
In~Section~\ref{sub:classical} we recall classical results on manifolds and graphs. We compare them with our Theorems~\ref{thm:normalizedglobal} and~\ref{thm:countingglobal}, which are special cases of our main result for the normalizing and counting measure, in Section~\ref{sub:globalmain}. The following Section~\ref{subsec:setup} introduces our general set-up and Section~\ref{section:mainresults} presents the general localized main result Theorem~\ref{thm:generaldeg1} for general measures. The strategy of the proof is discussed in Section~\ref{sec:strategy} and there the structure of the present work is discussed in detail.

\subsection{Classical results}\label{sub:classical}
Consider the heat equation
\[
\partial_t u =-\Delta u
\]
for a function $u=u(t,x)$, where $t>0$, $x\in \RR^n$, and $\Delta=-\sum_{i=1}^n\partial_i^2 $. Its fundamental solution is given by the Gauss-Weierstrass function
\[
f(t,x)=\frac{1}{\sqrt{4\pi t^n}}\exp\Bigl(-\frac{|x|^2}{4t}\Bigr),
\]
and the function $p_t(x,y)=f(t,x-y)$ is called the heat kernel. It constitutes the integral kernel of the heat semigroup $(\euler^{-t\Delta})_{t\geq 0}$ acting in $L^2(\RR^n)$.

Many spaces come with a natural notion of Laplace operator and an associated heat kernel. One might ask under which conditions on the space the heat kernel satisfies estimates comparable to the Gauss-Weierstrass function, and if these bounds can be geometrically characterized. As~mentioned above, there have been numerous works devoted to the geometric characterization of heat kernels in various settings. We will review some contributions related to the present work.

We start with the case where the underlying space is a Riemannian manifold.
Let $n>2$ and $M$ be a geodesically complete non-compact $n$-dimensional smooth Riemannian manifold with geodesic distance $d$ and volume measure $m$, and denote geodesic balls by $B_x(r)$ for $x\in M$ and $r>0$. Denote by $\Delta\geq 0$ the corresponding Laplace--Beltrami operator with heat semigroup $(\euler^{-t\Delta})_{t\geq 0}$ acting in $L^2(M,m)$, and minimal heat kernel $p$.

For the heat kernel $p$ of $M$, we~say that Gaussian upper bounds hold if there is a constant $C>0$ such that, for all $t>0$, $x,y\in M$, we~have
\begin{equation}\tag{$\rG_M$}\label{eq:GM}
p_t(x,y)\leq \frac{C\left(1+{d(x,y)^2}/{t}\right)^{\sfrac{n}{2}}}{\sqrt{m(B_x(\sqrt t))m(B_y(\sqrt t))}}\exp\Bigl(-\frac{d(x,y)^2}{4t}\Bigr).
\end{equation}
Short-time and long-distance behavior of the heat kernel are captured by the exponential.
Note that the constant $1/4$ appearing in the exponent is sharp compared to the Gauss-Weierstrass function and comes with the cost of the correction term $\left(1+{d(x,y)^2}/{t}\right)^{n/2}$.

A geometric condition ensuring Gaussian heat kernel bounds can be formulated in terms of Sobolev inequalities.
We say that $M$ satisfies local Sobolev inequalities if there is a constant $C>0$ such that, for all $x\in M$, $r>0$ and $u\in\cC_c^\infty(B_x(r))$, we~have
\begin{equation}\tag{$\rS_M$}\label{eq:SM}
C\frac{m(B_x(r))^{2/n}}{r^2}\|u\|_{\spfrac{2n}{n-2}}^2\leq \||\nabla u|\|_2^2+\frac{1}{r^2}\|u\|_2^2.
\end{equation}
Furthermore, the manifold $M$ satisfies the volume doubling property if there exists a constant $C>0$ such that, for all $x\in M$ and $0<r<R$, we~have
\begin{equation}\tag{$\rV_M$}\label{eq:VM}
\frac{m(B_x(R))}{m(B_x(r))}\leq C\Bigl(\frac{R}{r}\Bigr)^n.
\end{equation}

For instance, if $M$ has non-negative Ricci curvature, the properties \eqref{eq:GM}, \eqref{eq:SM} and \eqref{eq:VM} are simultaneously satisfied. In~general, one still has
the following characterization of Gaussian bounds for the heat kernel which can be extracted from {\cite{Grigoryan-94,BakryCLS-95}}, \cf \cite{SaloffC-01}:
\[
\eqref{eq:SM} \iff \eqref{eq:GM}\ \& \ \eqref{eq:VM}.
\]
The present paper focuses on Gaussian upper bounds of continuous-time heat kernels of discrete graphs. In~order to compare our results with the existing literature, we~briefly review existing results on this topic.

We consider locally finite connected infinite graphs with vertex set $X$, edge weights~$b$ and vertex measure $m$, see Section~\ref{subsec:setup} for definitions. There is a natural notion of Laplace operator and associated minimal heat kernel.

Analogous to the case of Riemannian manifolds, we~start with the heat kernel of a model space we want to compare our results with.
Let the canonical graph over the set of integers $\ZZ$ with standard edge weights be given, where $b(x,y)=1$ iff $|x-y|=1$, $x,y\in\ZZ$ and zero otherwise, and $m=1$. Denote by $d=|x-y|$ the combinatorial distance between $x,y\in\ZZ$.
An essentially exact estimate for $p$ was obtained in \cite[Th.\,3.5]{Pang-93}: there exists a constant $c>0$ such that
\[
\frac{c^{-1}}{\sqrt {t\vee d}}\exp\left(-2t\ \zeta\left(\sfrac{d}{2t}\right)\right)\leq p_t(x,y)\leq \frac{c}{\sqrt {t\vee d}}\exp\left(-2t\ \zeta\left(\sfrac{d}{2t}\right)\right)
\]
for $x,y\in\ZZ$, $d,t>0$,
where
\[
\zeta(x)
=
x\arsinh(x)
+1-\sqrt{x^2+1},\qquad x\geq 0.
\]
In particular, for $d>0$, we~have
\[
2t\ \zeta\left(\sfrac{d}{2t}\right)\sim \frac{d^2}{4t},\qquad t\to \infty.
\]
Here, $\sim$ means that left-hand side divided by right-hand side converges to one.

Next, we~consider a general graph $b$ over vertex set $X$ together with the normalizing measure, \ie $m=\deg=\sum_{y\in X}b(\cdot, y)$. Davies \cite{Davies-93} found that for any such graph, we~have with the combinatorial distance $d$, for all $t\ge 0$,
\[
p_t(x,y)\leq \frac{1}{\sqrt{\deg(x)\deg(y)}}\exp\bigl(-t\zeta\left(\sfrac{d(x,y)}{t}\right)\bigr).
\]
Note that the constant two in above estimate is missing due to a different scaling of the metric.

For the normalizing measure, a~characterization of heat kernel upper bounds in terms of Sobolev inequalities under volume growth assumptions for graphs can be found, e.g., in \cite{Barlow-book}. There, the following Gaussian bounds are considered: there is a constant $C>0$ such that, for all $x,y\in X$, we~have
\begin{equation}\tag{$\rG_U$}\label{eq:GU}
p_t(x,y)\leq \frac{C}{t^{n/2}}\exp\Bigl(-\frac{d(x,y)^2}{Ct}\Bigr),\qquad t\geq d(x,y).
\end{equation}
Note that the volume function is replaced by a uniform bound independent of the spatial variable. The off-diagonal bound holds only if $t\geq d(x,y)$, and moreover, instead of function $\zeta$, there is a constant multiple of $d(x,y)^2/t$.

For $n>2$, the characterization in \cite{Barlow-book} is based on the following global Sobolev inequality: there is a constant $C>0$ such that, for all $u\in \cC_c(X)$, we~have
\begin{equation}\tag{$\rS_U$}\label{eq:SU}
\|u\|_{\spfrac{2n}{n-2}}\leq C\||\nabla u|\|_2.
\end{equation}
If $n>2$ and the graph has Euclidean volume growth, \ie
\[
c^{-1}r^n\leq m(B_x(r))\leq c r^n,\qquad r\geq 1,\ x\in X,
\]
then we have
\[
\eqref{eq:GU} \iff \eqref{eq:SU}.
\]
The developed ideas employed to obtain upper bounds on the heat kernel have proved very fruitful in the study of continuous-time random walks in random environments, see, e.g., \cite{AndresDS-16,ChenKW-20,Bella-22} and the references therein. Explicitly, we~mention Barlow's work on Gaussian bounds on the heat kernel with normalizing measure on supercritical percolation clusters in $\ZZ^d$, \cite{Barlow-04}. A general feature of such random environments on graphs is that the involved conditions required to obtain heat kernel bounds only hold for sufficiently large balls.

Our main results are characterizations of upper heat kernel bounds on graphs on large scales involving $\zeta$ -- the sharp exponential used by Davies and Pang. In~the next section, we~discuss global results for the two important cases of normalizing and counting measure.

\subsection{Global main results for normalizing and counting measure}\label{sub:globalmain}
In this section, we~discuss two special cases of the main results.
Specifically, we~obtain versions of the above results neither assuming Euclidean volume growth nor boundedness of the degree. In~view of the discussion above, we~obtain these characterizations for sufficiently large balls and times in a quantitative manner.

We formulate our results in the spirit of the characterization for Riemannian manifolds described above. Given a graph {$ b $ over $ (X,m) $, see Section~\ref{subsec:setup} for details}, we~specify the meaning of Gaussian upper bounds, Sobolev inequalities and volume doubling before we state our results.

First, we~consider the case of a graph {$ b $} with {\it normalizing measure} $m=\deg$ and the combinatorial distance {$ d $. Furthermore, let $n>2$ be given.}

For $t_0\geq 0$,
the heat kernel $p$ satisfies Gaussian upper bounds if there is a constant $C>0$ such that for all $t\geq t_0$, $x,y\in X$ we have
\begin{equation}\tag{$\rG_{\norm}(t_0)$}\label{eq:Gnormt0}
p_t(x,y)\leq \frac{C\bigl(1+\sqrt{{d}(x,y)^2+t^2}-t\bigr)^{\sfrac{n}{2}}}{\sqrt{m(B_x(\sqrt t))m(B_y(\sqrt t))}}\exp\bigl(-t\zeta\left(\sfrac{{d}(x,y)}{t}\right)\bigr).
\end{equation}
Note that the involvement of the function $\zeta$ which is optimal on the integers leads to a change of the polynomial correction term. In~particular, the exponential term is different from the function known in the manifold case, but the polynomial correction term is smaller.

For $r_0\geq 0$, the graph satisfies a local Sobolev inequality if there is a constant $C>0$ such that for all $x\in {X}$, $r\geq r_0$, we have for all functions of $u$ supported in~$B_r(x)$
\begin{equation}\tag{$\rS_{\norm}(r_0)$}\label{eq:Snormr0}
C\,\frac{m(B_x(r))^{2/n}}{r^2}\|u\|_{\spfrac{2n}{n-2}}^2\leq \||\nabla u|\|_2^2+\frac{1}{r^2}\|u\|_2^2.
\end{equation}
Further, the graph satisfies the volume doubling property if there exists a constant $C>0$ such that, for all $x\in X$ and $0<r<R$, we~have
\begin{equation}\tag{$\rV_{\norm}$}\label{eq:Vnorm}
\frac{m(B_x(R))}{m(B_x(r))}\leq C\Bigl(\frac{R}{r}\Bigr)^n.
\end{equation}
Our main result for the normalizing measure is as follows.
\begin{theorem}[Normalizing measure, global version]\label{thm:normalizedglobal}Let $n>2$. If~there is $r_0>0$ such that \eqref{eq:Snormr0} holds, then there is $t_0=t_0(r_0)>0$ such that \eqref{eq:Gnormt0} and \eqref{eq:Vnorm} hold. Conversely, if there is $t_0>0$ such that \eqref{eq:Gnormt0} and \eqref{eq:Vnorm} hold, then there is $r_0=r_0(t_0)>0$ such that \eqref{eq:Snormr0} holds.
\end{theorem}

Our next main theorem concerns graphs with {\it counting measure} $m=1$. In~contrast to the normalizing measure, in presence of possibly unbounded vertex degree, it is essential to consider an intrinsic metric $\rho$ instead of the combinatorial distance~$ d $, \cf Section~\ref{subsec:setup}. For simplicity, we~assume now that $\rho$ is an intrinsic path metric with jump size less than one. For the implication \eqref{S}~$\Rightarrow$~ \eqref{G} \& \eqref{V}, it turns out that properties \eqref{G} and \eqref{V} need to be modified by an error involving the vertex degree. This error term is decreasing in time and radius. The implication \eqref{G}\ \&\ \eqref{V}~$\Rightarrow$~\eqref{S} requires the Sobolev dimension to be depending on the vertex degree inside the considered ball in order to obtain well-behaved Sobolev constants. To this end, we~employ a local regularity condition which naturally enters the characterization. Under mild assumptions, the variable Sobolev dimension is comparable to the doubling dimension.

Let $r_0,t_0\geq 0$ and $n>2$.
We say that the heat kernel $p$ satisfies Gaussian upper bounds if there is a constant $C>0$ such that for all $t\geq t_0$ and all $x,y\in X$ the heat kernel has the upper bound
\begin{multline}\tag{$\rG_{\rcount}(t_0)$}\label{eq:Gcountt0}
p_{t}(x,y)
\leq C(1\vee \deg(x)\vee \deg(y))
^{\sfrac{C}{\!\sqrt[ n]{t}}}\\
\cdot
\frac{\bigl(1+\sqrt{t^2+\rho(x,y)^2}-t\bigr)^{\sfrac{n}{2}}}
{\sqrt{m(B_{x}(\sqrt {t}))m(B_{y}(\sqrt {t}))}}\,
\euler^{-t\zeta(\rho(x,y)/t)}.
\end{multline}
These bounds will be characterized in terms of the following local Sobolev inequalities for a dimension function $n'\colon X\times [r_0,\infty)\to(2,\infty)$: there exists a {constant} $C\geq 1$ such that for all $x\in X$, $r\geq r_0$ and all functions of $u$ supported in $B_r(x)$, we have
\begin{equation}\tag{$\rS_{\rcount}(r_0)$}\label{eq:Scountr0}
C\,\frac{m(B_x(r))^{\sfrac{2}{n'_x(r)}}}{r^2}\,
\Vert u\Vert_{\spfrac{2n'_x(r)}{n'_x(r)-2}}^2
\leq \Vert \vert\nabla u\vert\Vert_2^2+\frac{1}{r^2}\Vert u\Vert_2^2.
\end{equation}
The {volume doubling} property is satisfied in $X$ if there is a constant $C>0$ such that for all $x\in X$ and $r_0\leq r_1\leq r_2$
\begin{equation}\tag{$\rV_{\rcount}(r_0)$}\label{eq:Vcountr0}
m(B_x(r_2))\leq C(1\vee\deg(x))
^{\sfrac{C}{\!\sqrt[ n]{r_1}}}\Bigl(\frac{r_2}{r_1}\Bigr)^{n} m(B_x(r_1)).
\end{equation}
The local regularity property is satisfied in $ X $ if there is a constant $C>0$ such that for all $x\in X$ and $r\geq r_0$, we~have
\begin{equation}\tag{L$(r_0)$}\label{eq:Lr0}
\frac{m(B_x(r))}{m(x)r^n}
\leq C\left(1\vee\deg(x)\right)^{\sfrac{n}2}.
\end{equation}
Our global main theorem in case of the counting measure is as follows.

\begin{theorem}[Counting measure, global version]\label{thm:countingglobal}
If $r_0>0$ and $n>2$ are constants and \eqref{eq:Scountr0} with dimension $n$ holds in $X$, then there exists \hbox{$t_0=t_0(r_0)>0$} such that \eqref{eq:Gcountt0}, \eqref{eq:Vcountr0} and \eqref{eq:Lr0} hold in $X$. Conversely, if $n{>2}$, $r_0,t_0>0$ are constants such that \eqref{eq:Gcountt0}, \eqref{eq:Vcountr0} and \eqref{eq:Lr0} hold in $X$, then there exists $r_1=r_1(r_0,t_0)>0$ such that $\rS_\rcount(r_1)$ holds with dimension function~$n'$~given~by
\[
n_x'(r)
=
n \Bigl(1+284n\frac{\ln(\ln r) }{\ln(r) } \Bigr)\Bigl(1+\frac{1}{2}\cdot\frac{\ln d_{x}(r)}{\ln r}\Bigr),
\]
where $d_x(r)=\max_{B_x(r)}(1\vee\deg)$. In~particular, if $\max_{B_x(r)}\deg\leq C r^k$ for some $C\geq 1$, $k\geq 0$ and all $r>0$, then
\[
n\leq \limsup_{r\to\infty}n'_x(r)\leq n\Bigl(1+\frac{k}{2}\Bigr).
\]
\end{theorem}

\begin{remark}
Observe that if $\max_{B_x(r)}\deg\leq C\ln r$ for some $C>0$ and all $r\geq r_0$, then $n'_x(r)\to n$ as $r\to\infty$. In~particular, if the vertex degree is bounded, we~recover the dimension $n$.
\end{remark}

\subsection{Set-up}\label{subsec:setup}

Let $X$ be an at most countable set and denote by $ \mathcal{C}(A) $ the real-valued functions $ f $ on $ X $ with support $ \mathrm{supp}\, f $ in $ A \subset X $ and by $ \mathcal{C}_{c}(X) $, we~denote the functions of compact support. {For a function $ f\in \mathcal{C}(A) $ which satisfies $ f\ge 0 $ or is absolutely summable, we~write $ \sum_{A}f=\sum_{x\in A}f(x) $.}
We extend a function
$m\colon X\to(0,\infty)$ to a measure on $X$ via \begin{align*}
m(A)=\sum_{A}m,
\end{align*} for $A\subset X$. {As usual, we~let
\[
\|f\|_p^p= {\sum_{X}m|f|^p},\qquad p\in[1,\infty),\qquad \|f\|_\infty=\sup_X|f|,
\]
and $\ell^p(X)=\{f\in \cC(X)\colon \|f\|_p<\infty\}$ for $p\in[1,\infty]$.}
We call a symmetric function $b\colon X\times X\to [0,\infty)$ such that $b(x,x)=0$ and
\[
\deg(x)=\sum_{y\in X}b(x,y)<\infty, \quad x\in X,
\]
a \emph{graph} over the measure space $(X,m)$. {We write $ x\sim y $ whenever $ b(x,y)>0 $ for $ x,y\in X $.} Furthermore, we~denote
\begin{align*}
\mathrm{Deg}(x)=\frac{1}{m(x)}\sum_{y\in X}b(x,y)=\frac{\deg(x)}{m(x)}.
\end{align*}

A graph is called \emph{locally finite} if the set $\{y\in X\colon b(x,y)>0\}$ is finite for any $x\in X$. The graph is called \emph{connected} if for any $x,y\in X$ there exists a finite sequence {$x=x_0\sim x_1\sim \ldots\sim x_n=y$ which we call a \emph{path} from $ x $ to $ y $.}

It is vital to use intrinsic metrics to deal with unbounded Laplacians on graphs, see \cite{Davies-93a,Folz-11,GrigoryanHuangMasamune-12,BauerKH-13, HaeselerKW-13, Folz-14b, Folz-14, BauerKW-15,Keller-15, HuangKS-20, KellerLW-21}.
An \emph{intrinsic metric} with respect to $b$ over $(X,m)$ is a non-trivial pseudo-metric $\rho\colon X\times X\to[0,\infty)$ such that
\[
\sum_{y\in X} b(x,y)\rho^2(x,y)\leq m(x),
\]
for all $ x\in X $.
A pseudo metric $ \rho $ is called a \emph{path metric} with {respect to the graph} if there is $ w: X\times X\to [0,\infty] $ such that for all $ x,y\in X $
\begin{align*}
\rho (x,y)=\inf_{x=x_{0}\sim \ldots\sim x_{n}=y}\sum_{j=1}^{n}w(x_{j-1},x_{j})
\end{align*}
{and $w(x,y)<\infty $ iff $ x\sim y $.}
Observe that the choice $ w(x,y)=(\Deg(x)\vee\Deg(y))^{-1/2} $ for $ x\sim y $ and $ w(x,y) =\infty$ otherwise yields an intrinsic path metric.

As usual, we~let
\[
B(r)= B_x(r)=\{y\in X: \rho(x,y)\leq r\},
\]
$r\geq 0 $, $x\in X$. Furthermore, we~will use the following notation for the space-time cylinder
\[
{Q(r_1,r_2)}=Q_x(r_1,r_2)=B_x(r_2)\times [r_1,r_2],
\]
for $ x\in X $ and $ r_1,r_2\in\RR$.

\begin{assumption}\label{assumption2} We assume that the intrinsic metric is a path metric for which the distance balls are compact and we have the \emph{finite jump size} condition
\[
S:=\sup\{\rho(x,y)\colon x,y\in X, b(x,y)>0\}<\infty.
\]
\end{assumption}

Observe that as a consequence our graphs are \emph{locally finite} and \emph{connected}. Indeed, local finiteness follows from finite balls and finite jump size, while connectedness follows from the fact that $ \rho $ is a path metric taking values in $ (0,\infty) $, \cf \cite{KellerLW-21}.

An important consequence of the assumptions above is that the metric space $ (X,\rho) $ is complete and geodesic, \ie for any two vertices $ x,y\in X $ there is a path $x= x_{0}\sim\ldots \sim x_{n}=y $ such that $ \rho(x,y)=\rho(x,x_{j})+\rho(x_{j},y) $ for all $ j=0,\ldots,n $, see \cite[Ch.\,11.2]{KellerLW-21} or \cite{KellerMuench}.

{For a locally finite graph, we~consider the {Laplace} operator
\[
\Delta f(x)= \frac{1}{m(x)}\sum_{y\in X}b(x,y)(f(x)-f(y)),\qquad f\in \mathcal{C}(X),\;x\in X.
\]
With slight abuse of notation, we~{also} denote by $ \Delta\geq0 $ the Friedrichs extension of the restriction of this operator to $ \mathcal{C}_{c}(X) $} in the Hilbert space $\ell^2(X,m)$. We denote by
\[
\Lambda:=\inf\spec(\Delta)
\]
the bottom of the spectrum of $\Delta$.

The minimal positive fundamental solution $p\colon [0,\infty)\times X\times X\to [0,\infty)$ of the heat equation
\[
\frac{d}{dt}u=-\Delta u \quad \text{on}\ [0,\infty)\times X
\]
is called the heat kernel of the graph. It can be seen via functional calculus that $ p $ is the kernel of the semigroup $(P_t)_{t\geq0}$, where
\begin{align*}
P_{t}=e^{-t\Delta},\qquad t\ge0.
\end{align*}
For $ x,y\in X $ and $ f\in\cC(X) $, {we let $ \nabla_{xy}f=f(x)-f(y) $ and}
\[
\vert \nabla f\vert(x):=\biggl(\frac{1}{m(x)}\sum_{y\in X}b(x,y) (\nabla_{xy}f)^2\biggr)^{\sfrac12}.
\]

In the following, we~introduce the conditions with which we will be dealing in this article. They generalize the well-established variants of {the notions discussed above} by introducing certain control functions.
\begin{definition}\label{def:sobvol} Let $B\subset X$, $R_2\geq R_1\geq 0$, as well as $n\colon X\times [R_1,R_2]\to (2,\infty)$, $\phi\colon X\times [R_1,R_2]\to [1,\infty)$, and ${\Psi}\colon X\times {(2,\infty)}\times [R_1,R_2]\to [1,\infty)$.
\begin{enumerateA}
\setcounter{enumi}{18}
\item\label{S}
The {\em Sobolev inequality} $\rS_\phi(n,R_1,R_2)$ holds in $B$, if for all $x\in B$, $r\in[R_1,R_2]$ {and $u\in\cC_{c}(B_x(r))$}, we~have
\begin{equation*}
\frac{m(B_x(r))^{\sfrac{2}{n_x(r)}}}{\phi_x(r)r^2}\,
\Vert u\Vert_{\spfrac{2n_x(r)}{n_x(r)-2}}^2
\leq \Vert \vert\nabla u\vert\Vert_2^2+\frac{1}{r^2}\Vert u\Vert_2^2.
\end{equation*}

\setcounter{enumi}{6}
\item\label{G} \emph{Gaussian upper bounds} $\rG_\Psi(n,R_1,R_2)$ are satisfied in $ B $ if for all $t\geq R_1^2$ and all $x,y\in B$ the heat kernel has the upper bound
\begin{multline*}
p_{t}(x,y)
\leq {\Psi_x(n_x(\sqrt \tau),\sqrt \tau)\Psi_y(n_y(\sqrt \tau),\sqrt \tau)}
\\
\cdot
\frac{\left(1\vee S^{-2}\left(\sqrt{t^2+\rho_{xy}^2S^2}-t\right)\right)^{\sfrac{n_{xy}(\sqrt \tau)}{2}}}
{\sqrt{m(B_{x}(\sqrt \tau))m(B_{y}(\sqrt \tau))}}\,
\euler^{-\Lambda (t-\tau)-(t/S^2)\zeta\left(\rho_{xy}S/t\right)},
\end{multline*}
where we set {$\tau:=t\wedge R_2^2$}, $ \rho_{xy}:=\rho(x,y) $, $ n_{xy}(t):=\frac12({n_x(t)+n_y(t)}) $
and recall
\begin{align*}
\zeta(x)=x\arsinh(x)+1-\sqrt{x^2+1}, \qquad x\geq 0.
\end{align*}

\setcounter{enumi}{21}
\item\label{V} The \emph{volume doubling} property $\rV_{{\Psi}}(n, R_1,R_2)$ is satisfied in $B$ if for all $x\in B$
\[
m(B_x(r_2))\leq {\Psi_x(n_x(r_2),r_1)}\Bigl(\frac{r_2}{r_1}\Bigr)^{n_x(r_2)} m(B_x(r_1)),\qquad R_1\leq r_1\leq r_2\leq R_2.
\]
Further, we~abbreviate $\rV_{{\Psi}}(n,R_2):=\mathrm{V}_{{\Psi}}(n,0,R_2)$, where we interpret $r/0=:\infty$ for $r\geq 0$.

\setcounter{enumi}{11}
\item\label{L}
The \emph{local regularity property} $\rL_\phi(n,R_1,R_2)$ is satisfied in $ B $ if for all $x\in B$ and $r\in[R_1,R_2]$, we~have
\[
\frac{m(B_x(r))}{m(x)(1\vee r)^{n_x(r)}}
\leq \phi_x(r)\left(1\vee\Deg(x)\right)^{\sfrac{n_x(r)}2}.
\]
\end{enumerateA}
\end{definition}
\begin{notation}
If {in the following} the functions $n,\phi,\Psi$ in the definitions above are chosen to be constant, we~will mention this explicitly.
\end{notation}

\begin{remark}[{Sharpness of $ \zeta $}]By the work of Davies and Pang, the function $\zeta$ in the case $S=1$ is sharp for the normalizing measure on the integers,
\cite{Davies-93, Pang-93}. Moreover, we~have for $r>0$
\[
t\zeta(r/t)\sim \frac{r^2}{2t},\qquad t\to \infty,
\]
where {$\sim$}
means that the left-hand side divided by the right-hand side converges to one.
\end{remark}

\begin{remark}[{The regularity property}] The local regularity property \eqref{L} can be interpreted as doubling property from balls with large radius to balls with very small radius. It should mainly be thought as $ m(B_{x}(r))/m(x)\leq C_x r^{n} $ which however fails to be equivalent for small radii $ r $. See \cite{BarlowChen-16} for a uniform version of \eqref{L}.
\end{remark}

\subsection{Localized main result for general measures}\label{section:mainresults}

{
In this section, we~present the main result which extends Theorems~\ref{thm:normalizedglobal} and \ref{thm:countingglobal} in two ways. First, we~allow for general measures and second, we~present a localized version of the characterization of heat kernel bounds. This means that \eqref{G}, \eqref{V}, and \eqref{S} now depend on the measure and are assumed to hold on annuli only.
Theorems~\ref{thm:countingglobal} as well as localized versions thereof follow from the general theorem below. In~order to formulate it, we~first need some preparations.
}

{Let $x\in X$ and}
a dimension function {$n\colon X\times [0,\infty)\to (2,\infty)$}, we~let the supremum over annuli {of radii} be given as
\begin{align*}
N_x(r)=\sup_{[r/4,r]} n_x.
\end{align*}
Furthermore, the volume doubling and {Gaussian} correction functions are given in terms of
\[
\Psi_x(N,r)=\left(1\vee\Deg(x)
\right)^{7N^2\sqrt[N]{S/r}}.
\]
Denote
the dimension function by
\[
n_x'(r)=
N_x(r) \Bigl(1+\frac{\ln(r'/(1\wedge r))}{\ln(r/r')} +{28
0 (1\vee S) N_x(r){\frac{\ln (1\vee r)}{\sqrt[ N_x(r)]{r'}}}}\Bigr)\Bigl(1+ \frac{1}{2}\frac{D_x(r)} {\ln \left(1\vee r\right)}\Bigr),
\]
where $D_x(r)=\sup_{B_x(r)}1\vee \Deg$ and
\[
r'= R_1\vee\Bigl( \frac{r}{4}\wedge \frac{1}{4}(\ln r)^{2\tilde N(r)}\Bigr),\qquad \tilde N(r)=\sup_{B_{x}(r)\times [R_1\vee (\ln r)/4,r]}n.
\]

The backbone of this article is the following new characterization of Gaussian upper heat kernel bounds on graphs.


\begin{theorem}[General locally regular case]\label{thm:generaldeg1}
Let
$R_2\geq R_1\geq {1024S}$, and $ B\subset X $.
\begin{enumeratei}
\item\label{thm:generaldeg1i}
If there is a constant $ \phi\ge 1 $ such that $\rS_\phi (n,R_1,R_2)$ holds in $B$, then {there exists a function $A=A(N,\phi,S)>0$} such that $\rG_{A\Psi}(N,4R_1,R_2)$, $\rV_{\!A\Psi}(N,R_1,R_2)$, and $\rL_{{A}}(N,R_1,R_2)$ hold in $B$.
\item\label{thm:generaldeg1ii}
If there is a function $A\!\geq\!1\mkern-2mu$ such that properties $\rG_{A\Psi}(n,R_1,R_2)$, $\rV_{\!A\Psi} (n,R_1,R_2)$, and $\rL_A(N,R_1,R_2)$ hold in $B=B_x(R_2)$, then there exists a positive function $\phi'=\allowbreak\phi'(N,A,S)>0$ such that $\rS_{\phi'}(n',4 R_1,R_2)$ holds in $x$.
\end{enumeratei}
\end{theorem}


\begin{remark}[Dimension function $n'$]
In Theorem~\ref{thm:countingglobal} we have seen that the dimension function $ n' $ converges to the doubling dimension if the vertex degree is bounded or if the vertex degree grows at most polynomially. In~contrast, here, additionally the input dimension $ n $ is a function. Assuming that $n$ is bounded then the same considerations can be applied to the dimension function $ n' $ as in Theorem~\ref{thm:countingglobal}. If~$n$ is unbounded, then one would naturally expect $n'$ to be unbounded as well. In~this case, one can still estimate $n' \leq C_\varepsilon N^{3+\varepsilon} D$ for all $\varepsilon>0$ and if $D$ grows at least logarithmically, then one can estimate $n' \leq C N^{2} D$.
\end{remark}

\begin{remark}[Functions $A$ and $\phi'$]Theorem~\ref{thm:generaldeg1} holds with
$A=(2\phi(1\vee S))^{54N^3}$ for part~\eqref{thm:generaldeg1i}. The statement in part \eqref{thm:generaldeg1ii} holds with Sobolev constant $\phi'$ given by the expression $\phi_x'(r) =2^{87+2 \spfrac{N_x(r)}{N_x(r)-2}}\Vert A\Eins_{Q}\Vert_\infty ^{20}$, where $Q=B_x(r)\times [(\ln r)_+/4,r]$. In~particular, if $A$ is a constant and $ n\ge n_{0}>2 $, then $\phi'$ can be chosen as a constant. So, for constant $A$, the only situation where $ \phi' $ may be unbounded is if one aims for a Sobolev inequality with dimension $n $ approaching $ 2 $.
\end{remark}

{\begin{remark}[Full volume doubling]
Indeed, one can deduce
\[
\rV_{\!A\Psi}(n,R_2)=V_{A\Psi}(n,0,R_2)
\]
in \eqref{thm:generaldeg1i} at the cost of a more sophisticated error function $\Psi$. We refer the interested reader to Theorem~\ref{thm:fulldoubling}.
\end{remark}}

\begin{remark}[The local regularity property] Theorem~\ref{thm:generaldeg1} is a consequence of the even more general Theorem~\ref{thm:general} which does not incorporate the local regularity property. However, the local regularity property \eqref{L} is a natural consequence of the Sobolev inequality, \cf Lemma~\ref{lemma:gammanormalized}. Thus, it is natural to include it into the characterization of the Sobolev inequality. Indeed, {in the case of the normalizing measure}, a~uniform version of \eqref{L} already appeared in \cite{BarlowChen-16}.
\end{remark}

The results of this paper are applied to specific examples including the Laguerre operator, \cf \cite{Kostenko-21} in an upcoming paper \cite{KellerKNR-24}. There we use isoperimetric estimates to infer Sobolev inequalities to employ the heat kernel bounds derived here.

\subsection{Strategy of the proofs}\label{sec:strategy}

Theorems~\ref{thm:normalizedglobal}, \ref{thm:countingglobal} and \ref{thm:generaldeg1} are special cases of more general technical results. These are summarized in Theorem~\ref{thm:general} which does not involve a local regularity property. Roughly speaking we want to show an ``equivalence'' of \eqref{G} \& \eqref{V} \& \eqref{L}
and \eqref{S}. Below we will discuss the overall strategy and the necessary adaptions for the special cases.

Section~\ref{section:doubling} is dedicated to show how \eqref{S} implies \eqref{V} \& \eqref{L}. To deduce \eqref{V} from \eqref{S} we follow the strategy of \cite{SaloffC-01}. Inserting appropriate cut-off functions into \eqref{S} gives an intermediate step of an iteration procedure. The number of steps of this iteration procedure is restricted by the jump size. The error term for the volume doubling -- above named $ {\Psi} $ -- which emerges from the iteration depends now on {the vertex degree of the center of} the initial ball, \cf Theorem~\ref{thm:adapteddoubling}. As~discussed above $\Psi $ can be controlled for large radii. Furthermore, \eqref{L} is an immediate consequence of \eqref{S} by plugging in characteristic functions of vertices, \cf Lemma~\ref{lemma:gammanormalized}.

In Section~\ref{section:SG} we show how to derive \eqref{G} from \eqref{S} after we already obtained \eqref{V}. The considerations of this section are a variation of the strategy developed in \cite{KellerRose-22a}. This strategy is a combination of truncated Moser iteration and Davies' method. The significant difference to \cite{KellerRose-22a} is that our new error terms to estimate $ p_{t}(x,y) $ in \eqref{G} now only depend on the vertex degree of $ x $ and $ y $ rather than means in large balls about these vertices of radius $ \sqrt{t} $, \cf Theorem~\ref{thm:main1sturm}.

A strategy to derive \eqref{S} from \eqref{G} \& \eqref{V} goes back to \cite{Varopoulos-85}, \cf \cite{SaloffC-01} and infers a weak Sobolev inequality from uniform diagonal heat kernel bounds. A challenge in the discrete setting is that this strategy only gives a weak Sobolev inequality involving the uniform norm instead of the $ 1 $-norm.
This issue requires the incorporation of the trivial $\ell^\infty$-$\ell^1$-embedding on balls which yields an unpleasant error in terms of reciprocals of measures within the Sobolev constant. We then conclude \eqref{S} from the weak Sobolev inequality by a version of the proofs in \cite{Barlow-book, Delmotte-97}. The unpleasant error is still good enough in very bounded situations, \cf Theorem~\ref{thm:heattosobolevgeneral}. However, allowing for more unboundedness we need a new idea which is the choice of a variable dimension. This choice mitigates the possibly unbounded error terms for large radii, \cf Theorem~\ref{thm:heattosobolevgeneralimproving} which still includes a free parameter which is chosen later to prove Theorems \ref{thm:countingglobal} and~\ref{thm:generaldeg1}.

In Section~\ref{section:remaining} we then put Theorem~\ref{thm:adapteddoubling},~\ref{thm:main1sturm} and~\ref{thm:heattosobolevgeneralimproving} together to conclude the most general result Theorem~\ref{thm:general}. The proof for the normalizing measure, Theorem~\ref{thm:equivnormalized}, then employs a corollary of Theorem~\ref{thm:adapteddoubling},~\ref{thm:main1sturm} and~\ref{thm:heattosobolevgeneral}. For Theorem~\ref{thm:generaldeg1}, the strategy to conclude \eqref{G} \& \eqref{V} \& \eqref{L}
from \eqref{S} is similar and only invokes Lemma~\ref{lemma:gammanormalized} additionally. For the ``reverse direction'' we then have to choose the free parameter in Theorem~\ref{thm:heattosobolevgeneralimproving}. This allows us to show \eqref{S} with a variable dimension. Here, \eqref{L} yields that all error terms only depend on the vertex degrees and not on reciprocals of measures.

\subsubsection*{Acknowledgements}
The authors thank Tomasz Grzywny for an illuminating hint which improved the manuscript and the anonymous referee for the encouragement to further simplify the presentation.

\section{Sobolev implies volume doubling and local regularity}\label{section:doubling}

In this section we obtain volume doubling properties in balls where Sobolev inequalities are satisfied. The general strategy is based on \cite{SaloffC-01}. We show how the variable dimension of the Sobolev inequalities {translates into} the doubling dimension. As~another new feature, the influence of the local geometry on the doubling condition of the graph is made explicit. Finally, we~discuss how to derive the local regularity property from the Sobolev inequality.

Recall that $\rS_\phi(n,r,r)$ means that we have a Sobolev inequality with Sobolev constant $ \phi $ for the radius $ r $. We first show that the Sobolev inequality remains true when increasing the dimension $ n $. This is common knowledge but included for the convenience of the reader.

\begin{lemma}\label{lem:sobdim} Let $r>0$, $n>2$, $\phi>0$, $x\in X$ and assume $\rS_\phi(n, r,r)$ in $x$. Then $\rS_\phi(N,r,r)$ holds in $x$ for all $N\geq n$.
\end{lemma}
\begin{proof}
Let $f\in\cC_{c}(B(r))$. From Hölder's inequality with exponents
\[
p=\frac{n}{n-2}\frac{N-2}{N}\quad\text{and}\quad q^{-1}=\Bigl(\frac{n}{n-2}\frac{N-2}{N}-1\Bigr)\frac{N}{N-2}\frac{n-2}{n},
\]
\ie $ \sfrac1p+\sfrac1q=1 $, we get
\begin{align*}
\frac{m(B(r))^{\sfrac{2}{N}}}{r^2}&\,\Vert f\Vert_{\spfrac{2N}{N-2}}^2 =\frac{m(B(r))^{\sfrac{2}{N}}}{r^2}\Bigl(\sum_{X} m |f|^{\spfrac{2N}{N-2}}\cdot 1\Bigr)^{\psfrac{N-2}{N}}
\\
&\leq
\frac{m(B(r))^{\sfrac{2}{N}}}{r^2}\biggl(\Bigl(\sum_{X} m |f|^{p\spfrac{2N}{N-2}}\Bigr)^{\sfrac{1}{p}}m(B(r))^{\sfrac{1}{q}}\biggr)^{\psfrac{N-2}{N}}\\
&=\frac{m(B(r))^{\sfrac{2}{n}}}{r^2}\,\Vert f\Vert_{\spfrac{2n}{n-2}}^2.
\end{align*}
Since the right-hand side of $\rS_{\phi}(N,r,r)$ does not depend on $ N $, this yields the claim.
\end{proof}

Next, we~derive a Nash inequality from a Sobolev type inequality which is folklore {but} included for the readers convenience.

\begin{lemma}[Sobolev to Nash]\label{lem:StoN}
Let $r>0$, $n>2$, {$ C>0 $}. Assume that for $f\in \cC({X})$ we have
\[
\Vert f\Vert_{\spfrac{2n}{n-2}}^2\leq C\Bigl(\Vert |\nabla f|\Vert_2^2+\frac{1}{r^2}\Vert f\Vert_2^2\Bigr).
\]
Then, we~have
\[
\Vert f\Vert_2^{2+4/n}\leq C \Bigl(\Vert |\nabla f|\Vert_2^2+\frac{1}{r^2}\Vert f\Vert_2^2\Bigr)\Vert f\Vert_1^{4/n}.
\]
\end{lemma}

\begin{proof}
Recall the Lyapunov inequality, which is a consequence of the Hölder inequality: For $1\leq p_0,p_1<\infty$, $\theta\in(0,1)$ and $p\in [1,\infty)$ such that
\[
\frac{1}{p}=\frac{1-\theta}{p_0}+\frac{\theta}{p_1},
\]
we have for all $f\in\cC_c(X)$
\[
\Vert f\Vert_p\leq \Vert f\Vert_{p_0}^{1-\theta}\Vert f\Vert_{p_1}^\theta.
\]
If we choose $p=2$, $p_0=\spfrac{2n}{n-2}$, and $p_1=1$, then $\theta=\spfrac{2}{n+2}$ and $1-\theta=\spfrac{n}{n+2}$. The claim follows from applying the Sobolev inequality to the right-hand side of the resulting inequality.
\end{proof}

The following lemma shows that a Sobolev type inequality with dimension $ n $ yields a lower bound on $ m(B_{x}(r))/r^n $ in an interval.

\begin{lemma}[Non-collapsing]\label{lem:asnoncoll}Let $x\in X$, $R {\ge 4S}$,
$n>2$ and $C>0$ constants. Assume that for all $f\in \cC_{c}(B_x(R))$ we have
\[
\Vert f\Vert_{\spfrac{2n}{n-2}}^2\leq C\Bigl(\Vert| \nabla f|\Vert_2^2+\frac{1}{R^2}\Vert f\Vert_2^2\Bigr).
\]
Then we have for all $r\in[2S,R]$
\[
2^{-{3}n^2}\Bigl(\frac{1}{C}\Bigr)^{\sfrac{n}{2}}
\Bigl[1\wedge C^{\sfrac{n}{2}}\frac{m(x)}{r^{n}}\Bigr]^{2^{2}\sqrt[n]{{S}/{r}}}
\leq \frac{m(B_x(r))}{r^{n}}.
\]
\end{lemma}

\begin{proof}
Lemma~\ref{lem:StoN} yields for all $f\in \cC_{c}(B_x(R))$ the Nash inequality
\[
\Vert f\Vert_2^{2+4/n}\leq C \Bigl(\Vert |\nabla f|\Vert_2^2+\frac{1}{R^2}\Vert f\Vert_2^2\Bigr)\Vert f\Vert_1^{4/n}.
\]
We apply this to special cut-off functions.
For {$r\in[2S,R]$}, choose
\[
f_r(y):=
\begin{cases}
{r}-\rho(x,y),& y\in B_{x}(r/2),\\
\bigl(\sfrac{3r}2-2\rho(x,y)\bigr)_+,&
\mbox{else},
\end{cases}
\]
where $a_+=a\vee0$ for $a\in \mathbb{R}$. Then, $f_r$ satisfies {$\supp f_r\subset B_x(3r/4)\subset B_x(r)$},
\[
\Vert f_r\Vert_1\leq {r} m(r), \quad
\quad \frac{r}{2}\,m(r/2)^{1/2}\leq \Vert f_r\Vert_2\leq {r} m(r)^{1/2},
\]
where we used the notation $ m(r)=m(B_{x}(r)) $.
Since {$ r\ge 4S $}, we~have
$\supp\vert \nabla f_r\vert\subset B_x(3r/4+S)\subset B_x(r)$. Thus, since $\vert \nabla f_r\vert\leq {2}$,
\[
\Vert \vert \nabla f_r\vert \Vert_2
\leq {2} m(r)^{1/2}.
\]
Therefore, the Nash inequality applied to $f_r$ yields
\[
\Bigl(\frac{r}2 m({r}/2)^{\sfrac12}\Bigr)^{2+\sfrac4n}
\leq
C \Bigl(4m(r)+\frac{({r})^2}{R^2}m({r})\Bigr)
{r}^{\sfrac4n}\left( m({r})\right)^{\sfrac4n}
\leq 8 C \ m(r)^{1+\sfrac4n}
r^{\sfrac4n},
\]
where we used {$ r\le R $} and the monotonicity of the measure in the last line. If~we
put
\[
\alpha=1+\frac2n,\quad \beta=1+\frac4n,\quad\text{and}\quad q=\frac\alpha\beta=\frac{n+2}{n+4},
\]
this is equivalent to
\[
\Bigl(\frac{r^2}{2^{2\alpha+3}C}\Bigr)^{\sfrac{1}{\beta}} m({r}/2)^{q}
\leq m(r).
\]
Iterating the above inequality we obtain
\[
\Bigl(\frac{r^2}{2^{2\alpha+3}C}\Bigr)^{\Psfrac{1}{\beta}\sum_{i=0}^{k-1}q^i}\Bigl(\frac12\Bigr)^{\Psfrac{2}{\beta}\sum_{i=1}^{k-1}iq^i}
m({r}/{2^k})^{q^k}
\leq m(r).
\]
This iteration procedure yields a non-trivial lower bound on $m(r)$ as along as we have $r/2^{k-1}\geq 4S$, \ie if we choose
\[
k\leq
\Bigl\lfloor \log_2\frac{r}{2 S}\Bigr\rfloor =:\eta(r).
\]
We are left with the sums in the exponents and start with the calculation of the first from the left. Set $\theta:=q^{\eta}$. Clearly, geometric summation gives
\[
\sum_{i=0}^{{\eta}-1}q^i=\frac{1-q^{{\eta}}}{1-q}=\frac{1-{{\theta}}}{1-q}.
\]
Next, since $q\in(0,1)$ we have
\[
\sum_{{i}=1}^{k-1}iq^{i}\leq \sum_{i=0}^\infty iq^i=\frac{q}{(q-1)^2}.
\]
Hence, using $m(r)=m(B_x(r))\geq m(x)$ for all $r\geq 0$,
{\begin{align*}
m(B_x(r))&\ge\Bigl(\frac{r^2}{2^{2\alpha+3} C }\Bigr)^{\Psfrac{1}{\beta}\pspfrac{1- {\theta}}{1-q}}
\Bigl({\frac 12}\Bigr)^{\Psfrac{2}{\beta}\sfrac{q}{(1-q)^2}}
m(x)^{{\theta}}
\\
&=\Bigl(\frac{r^2}{C }\Bigr)^{\Psfrac{n}{2}(1- {\theta})}
{\Bigl(\frac12\Bigr)^{(\Psfrac52n+2)(1-\theta)+\sfrac{n^2}{2}+1}}
m(x)^{{\theta}},
\end{align*}}
where, by noting $q=\alpha/\beta$, $\alpha=1+2/n$ and $\beta=1+4/n$, we~used the trivial identities $\beta(1-q)=\beta(1-\alpha/\beta)=\beta-\alpha=2/n$ and
$q/(1-q)=\alpha/(\beta-\alpha)=n/2+1$.
Since $ {\theta}=q^{{\eta}}\in(0,1)$, we~get for the exponent of $1/2$ using $ n>2 $
\[
\Bigl(\frac52n+2\Bigr)(1-\theta)+\frac{n^2}{2}+n
\leq \frac{n^2}{2}+\frac{7}{2}n+2 \leq 3 n^2.
\]
This leads to the estimate
\[
2^{-{3}n^2}\Bigl(\frac{1}{C}\Bigr)^{\sfrac{n}{2}}
\Phi^{\theta }
\leq \frac{m(B_x(r))}{r^{n}}.
\]
where $\Phi= C^{{n}/{2}}{m(x)}/{r^{n}} $,
$ \theta =q^{\eta(r)}$, $q=(n+2)/(n+4)\in(0,1)$, and {$ \eta(r)=\lfloor \log_2 (r/2S) \rfloor$}. In~order to obtain the claim, we~estimate $\Phi\geq 1\wedge \Phi$ and bound $\theta$ from above. From $\lfloor x\rfloor\geq x-1$ for all $x\in\RR$ we infer
\[
\eta\geq \log_2{r/2S}-1\geq \log_2{r/4S}.
\]
{By using the elementary inequality $ \ln (1+x)\ge 2x/(2+x) $ for $ x\ge 0 $, we~obtain the inequality
$ n\ln ((n+4)/(n+2))\ge 2n/(n+3) $. Furthermore, $ 2n/(n+3)\ge \ln 2 $ for $ n>2 $. Thus, we~have
\[
\log_2({1/q})=\log_2({(n+4)/(n+2)})\ge 1/n.
\]
}
Therefore, since $S/r\leq1$,
\[
\theta=q^{\eta}\leq {q}^{\log_2(\sfrac{r}{4S})}=4^{\log_2({1/q})}\Bigl(\frac{S}{r}\Bigr)^{\log_2({1/q})}\leq 2^2\Bigl(\frac{S}{r}\Bigr)^{\sfrac{1}{n}}.
\]
Applying this estimate leads to the claim.
\end{proof}

The following lemma essentially appears already in \cite[Lem.\,6.2]{KellerRose-22a} for the normalizing measure.

\begin{lemma}[Local regularity]
\label{lemma:gammanormalized} Let $\phi\geq 1$, $n>2$ be constants, $r \ge 0$, and $o\in X$. If~$\rS_{\phi}(n,r,r)$ holds in $ o $, we~have for all $x\in B_o(r)$
\[
\frac{m(B_o(r))}{m(x)}
\leq {3}^{\sfrac{n}{2}}\phi^{\sfrac{n}{2}}(1\vee r)^n\left(1\vee \Deg(x)\right)^{\sfrac{n}{2}}.
\]
In particular, if $R_2\geq R_1\geq 0$, $B\subset X$ and $\rS_\phi(n,R_1,R_2)$ holds in $B$, then property $\rL_{\tilde\phi}(n,R_1,R_2)$ holds in $B$ with $\tilde \phi={3}^{n/2}\phi^{n/2}$.
\end{lemma}
\begin{proof}
This follows directly by applying $\rS_{\phi}(n,r,r)$ to $ u=1_{x} $ which is supported in $ B_{o}(r) $:
\begin{align*}
m(B_o(r))^{\sfrac{2}{n}} m(x)^{\psfrac{n-2}{n}}&\leq \phi\left({2} r^2\Deg(x)m(x) +m(x)\right)\\
&\leq {3} (1\vee r^2)(1\vee \Deg(x))m(x),
\end{align*}
which settles the claim.
\end{proof}

\begin{theorem}[Volume doubling]\label{thm:adapteddoubling} Let $x\in X$, $R_2\geq R_1\geq 4S$, $r\in[S,R_2/2]$, a~function $n:[R_{1},R_{2}]\to (2,\infty)$ and a constant $\phi\geq 1$. If~$\rS_{\phi}(n,R_1,R_2)$ is satisfied in $x\in X$, then we have $\rV_{\!A\Phi}(n,R_1,R_2)$ in $x$, \ie
\[
\frac{m(B_x(R))}{R^{n(R)}}
\leq {A(n(R))} \Phi_x(n(R),r)\ \frac{m(B_x(r))}{r^{n(R)}}, \quad R_1\leq r\leq R\leq R_2,
\]
where
\[
\Phi_x(n,r)=\left(1\vee\Deg(x)\right)^{2n\sqrt[n]{S/r}},\quad\mbox{and}\quad A(n)=2^{8n^2}\phi^{\sfrac{n}{2}}(1\vee S)^{4n}.
\]
Moreover, we~obtain for all $r\in[2R_1,R_2]$ the doubling property
\begin{align*}
m(B_x(r))\leq {A'(n(r))} \Phi_x(n(r),r/2) m(B_{x}(r/2))
\end{align*}
with $ A'(n)=2^{9n^2}\phi^{\sfrac{n}{2}}(1\vee S)^{4n} $.
\end{theorem}

\begin{proof}
Let $R\in[R_1,R_2]$, $n:=n(R)$, and $w:=2^{-2}\sqrt[n]{r/S}$. The non-collapsing lemma, Lemma~\ref{lem:asnoncoll}, yields with $C=\phi\sfrac{R^{2}}{m(B(R))^{\sfrac{2}{n}}}$ for all $r\in [ 4S,R]$
\[
\frac{1}{2^{{3}n^2}\phi^{\sfrac{n}{2}}}\frac{m(B_x(R))}{R^{n}}
\left[1\wedge \phi^{\sfrac{n}{2}}\frac{R^{n}}{m(B_x(R))}\frac{m(x)}{r^{n}}\right] ^{\sfrac{1}{w}}
\leq \frac{m(B_x(r))}{r^{n}}.
\]
By Lemma~\ref{lemma:gammanormalized}, we~have $\rL_{\tilde\phi}(n,R,R)$ in $x$ with $\tilde \phi={3}^{n{/2}}\phi^{n{/2}}$, \ie we have the estimate ${{m(B_x(R))}/{m(x)}
\leq {3}^{n{/2}}\phi^{n/2}(1\vee R^n)\left(1\vee\Deg(x)\right)^{{n}/{2}}}$.
Division by the first and third factor and applying the bound above yields
\begin{align*}
\frac{m(B_x(R))}{R^{n}}
&\leq
2^{{3}n^2}
\phi^{\sfrac{n}{2}}
\left[1\vee \left(\frac{r}{R}\right)^{n} {\frac{m(B_x(R))}{\phi^{\sfrac{n}{2}}m(x)}}\right]^{\sfrac{1}{w}}\frac{m(B_x(r))}{r^{n}}
\\
&\leq
2^{{3}n^2}
\phi^{\sfrac{n}{2}}
\left[1\vee \frac{{r}^{n}}{R^{n}}
{3}^{n{/2}}(1\vee R^n)\left(1\vee\Deg(x)\right)^{\sfrac{n}{2}}
\right]^{\sfrac{1}{w}}\frac{m(B_x(r))}{r^{n}}
\\
&\leq
2^{{5}n^2}
\phi^{\sfrac{n}{2}} (1\vee r^n)^{\sfrac{1}{w}}
\left(1\vee\Deg(x)\right)
^{\sfrac{n}{2w}}\frac{m(B_x(r))}{r^{n}}.
\end{align*}
since $ n/w= 2^2n\sqrt[n]{S/r} \leq 2n^2$ and, thus, $ {3}^{n/{2}w} \leq 2^{2n^2}$.
Now, note
\[
r^{\frac{n}{w}}=S^{2^2n\sqrt[n]{S/r}}h(r/S)^{2^2n}
\]
with
\[
h(x)=x^{\sqrt[n]{1/x}}=\exp\Bigl(\frac{\ln x}{x^{\sfrac{1}{n}}}\Bigr).
\]
Differentiating with respect to $ x $ yields that the maximum is attained at $ x= \euler^{n}$ and hence we can estimate $ h(r/S)^{2^2n}\leq e^{2^2n^2/e}\leq 2^{3n^2} $ since $ e^{{4/e}}\leq 2^3 $. Now, use $\phi\geq 1$ to obtain the first claim.
The second claim is an application of the first one.
Just observe using $2<n$
\[
\frac
{m(B_x(r))}
{m(B_x(r/2))}\leq
2^{9n^2}\phi^{\sfrac{n}{2}}(1\vee\Deg(x))^{2n\sqrt[n]{2S/r}}.
\]
This finishes the proof.
\end{proof}

Given $0\leq R_1\leq R_2$, we~extend the function
$\Psi\colon X\times (2,\infty)\times [R_1,R_2]\to[1,\infty)$ to the function $\Psi=\Psi_{R_1}\colon X\times (2,\infty)\times(0,R_2]\to[1,\infty)$ by
\[
\Psi_x(n,r)
=\left(1\vee\Deg(x)\right)^{{2}n\sqrt[n]{S/(r\vee R_1)}}
\cdot\begin{cases}
\bigl((1\vee R_1^{2})(1\vee\Deg(x))\bigr)^{\sfrac{n}{2}},& r\leq R_1,
\\
1,& \text{else.}
\end{cases}
\]
With this modification of the function {$\Psi$}, we~obtain the following extension of volume doubling for small radii.
\begin{lemma}\label{lem:fulldoubling}
Let $x\in X$, $R_2\geq R_1\geq 0$, and functions $n:[R_{1},R_{2}]\to (2,\infty)$ and {$ C_1\colon (2,\infty)\to[1,\infty)$, $C_2\colon [R_1,R_2]\to[1,\infty) $}. If~$\rV_{C_1\Psi}(n,R_1,R_2) $ and {$\rL_{C_2}(n,R_1,R_1)$} hold in $x$, then we have $\rV_{C_3\Psi}(n,R_2)$ in $x$, where
\[
C_3(n(r))=(2C_2(R_1))^{\sfrac{n(R_1)}2}
\begin{cases}
1,& r\in[0,R_1],
\\
C_1(n(r)),& r\in(R_1,R_2],
\end{cases}
\]
where we extended $ n(r)=n(R_{1}) $ for $ r\in [0,R_{1}) $.
\end{lemma}
\begin{proof}
We have to check two cases $0<r_1\leq r_2\leq R_1$ and $0< r_1\leq R_1\leq r_2\leq R_2$ and start with the case $0\leq r_1\leq r_2\leq R_1$.
We infer from {$\rL_{C_2}(n,R_1,R_1)$} the estimate
\[
m(B(R_1))/{m(x)}\leq (2C_2(R_1)(1\vee R_1^2)(1\vee \Deg(x)))^{\sfrac{n(R_1)}{2}}.
\]
Therefore, for all $r_1,r_2$ such that $0< r_1\leq r_2\leq R_1$, we~have {since $C_1\ge 1$}
\begin{align*}
\frac{m(B(r_2))}{m(B(r_1))}\leq \frac{m(B(R_1))}{m(x)}&\leq
(2C_2(R_1)(1\vee R_{1}^{2})(1\vee\Deg(x)))^{\sfrac{n(R_1)}{2}}\\
&\leq C_3(r_2)\Psi_{x}(n(r_2),r_1).
\end{align*}
Further, if $0<r_1\leq R_1\leq r_2\leq R_2$, we~use $\rV_{C_1{\Psi}}(n,R_1,R_2)$ in $x$ and the second inequality in the above displayed formula for the choice $r_2=R_1$ as well as $r_1\leq R_1$ to obtain
\begin{align*}
\frac{m(B(r_2))}{m(B(r_1))}
&\leq
C_{1}(n(r_2)){\Psi}(n(r_2),R_1)\Bigl(\frac{r_2}{R_1}\Bigr)^{n(r_2)}\frac{m(B(R_1))}{m(B(r_1))}
\\
&\leq
C_{1}(n(r_2))\Psi(n(r_2),R_1)[2C_2(R_1)(1\vee R_1^2)(1\vee\Deg(x))]^{\sfrac{n(R_1)}{2}}\Bigl(\frac{r_2}{r_1}\Bigr)^{n(r_2)}.
\end{align*}
Thus, the claim follows.
\end{proof}
The above considerations thus lead to the following extension of Theorem~\ref{thm:adapteddoubling}.
\begin{theorem}[Full doubling]\label{thm:fulldoubling}
Let
\[
x\in X,\quad R_2\geq R_1\geq 4S,\quad r\in[S,R_2/2],\quad n:[R_{1},R_{2}]\to (2,\infty),
\]
and a constant $\phi\geq1$ be given. If~$\rS_\phi(n,R_1,R_2)$ holds in $x$, then $\rV_{\!A\Psi}(n,R_2) $ holds in~$x$, where
\[
\Psi_x(n,r)
=\left(1\vee\Deg(x)\right)^{{2}n\sqrt[n]{S/(r\vee R_1)}}
\cdot\begin{cases}
((1\vee R_1^{2})(1\vee\Deg(x)))^{\sfrac{n}{2}},& r\leq R_1,
\\
1,& \text{else,}
\end{cases}
\]
and
\[
A(n(r))=(2\phi)^{{n(R_1)^2} }2^{8n(r)^2}\phi^{\sfrac{n(r)}2}(1\vee S)^{4n(r)},
\]
where we extended $ n(r)=n(R_{1}) $ for $ r\in [0,R_{1}) $.
\end{theorem}
\begin{proof}
By Theorem~\ref{thm:adapteddoubling}, we~infer {$\rV_{\!\tilde A\Psi}(n,R_1,R_2) $ holds in $x$, where $\tilde A$ can be chosen as $\tilde A(n)=2^{{8}n^2}\phi^{n{/2}}(1\vee S)^{{4}n}$}. Moreover, Lemma~\ref{lemma:gammanormalized} yields, in particular, $\rL_{\tilde \phi}(n,R_1,R_1)$ in $x$ with $\tilde\phi={3}^{n(R_{1})/2}\phi^{n(R_{1})/2}$. Hence, Lemma~\ref{lem:fulldoubling} yields the claim with the constant
\begin{align*}
A'(n(r))&=2^{\sfrac{n(R_{1})}{2}}\,\tilde\phi^{\sfrac{n(R_{1})}{2}}\tilde A(n(r))\\
&=2^{\sfrac{n(R_1)}{2}}\,3^{\sfrac{n(R_1)^2}{4}}\,2^{8n(r)^2}\phi^{\sfrac{n(R_1)^2}{4}+\sfrac{n(r)}{2}}(1\vee S)^{{4}n(r)}\leq A(n(r))
\end{align*}
for $r\in[R_1,R_2]$. Analogously,
Lemma~\ref{lem:fulldoubling} yields the constant
\[
A'(n(r)) = 2^{\sfrac{n(R_{1})}{2}}\,\tilde\phi^{\sfrac{n(R_{1})}{2}}= 2^{\sfrac{n(R_1)}{2}+\sfrac{n(R_1)^2}{4}}\, \phi^{\sfrac{n(R_1)^2}{4}}\leq A(n(r))
\]
for $ r\leq R_{1} $. This proves the claim.
\end{proof}

Observe that for the normalizing measure, the combinatorial graph distance is an intrinsic metric. For this metric, the jump size is always $ 1 $ which explains the uniform lower bound on the smallest radius $ R_{1} $ in the next corollary.

\begin{corollary}[Full doubling of normalized measure]\label{cor:adapteddoubling} Let $m=\deg$ and $ \rho $ be the combinatorial graph distance. Let $\phi\geq 1$, $ n>2 $ be constants and $R_2\geq R_1\geq 4$. If~$\rS_{\phi}(n, R_1,R_2)$ holds in $x\in X$, then we have $\rV_{{\Psi}}(n, R_2)$ in $x$,
where
\[
\Psi = 2^{9n^2}\phi^{n^2} R_1^n.
\]
\end{corollary}
\begin{proof}
From Theorem~\ref{thm:fulldoubling}, we~obtain the property $\rV_{\!A{\tilde \Psi}}(n,R_{2}) $ in $x$ with the constant $ A=2^{{n^2}+8n^2}\phi^{n^2/2+n/2}(1\vee S)^{4n}\leq 2^{9n^2}\phi^{n^2} $ since $ S=1 $.
Since $ \Deg =1 $, we~have $ \tilde \Psi(n,r)=1 $ for $ r\ge R_1 $ and $ \tilde\Psi(n,r) =(1\vee R_{1})^{n}=R_{1}^{n}$ since $ R_{1}\ge 4 $. This settles the claim.
\end{proof}

\section{Sobolev implies Gaussian bounds}\label{section:SG}

In this paragraph we adjust and expand some of the techniques from \cite{KellerRose-22a}. In~order to obtain estimates of solutions of the heat equation, we~investigate properties of solutions of the \emph{$\omega$-heat equation}
\[
\frac{d}{dt} v_t= -\Delta_\omega v_t,
\]
where $\Delta_\omega:= \euler^\omega\Delta \euler^{-\omega}$ is a sandwiched Laplacian for $\omega\in\ell^\infty(X)$.
The following results provide an $\ell^2$-mean value inequality for non-negative solutions of the $\omega$-heat equation. The displacement of the solutions with respect to the heat equation is measured in terms of the function
\[
h(\omega)=\sup_{x\in X}\frac{1}{m(x)}\sum_{y\in X} b(x,y)\vert \nabla_{xy}\euler^\omega\nabla_{xy}\euler^{-\omega}\vert.
\]
The semigroup $P_t^\omega:=\euler^\omega P_t\euler^{-\omega}$, $ t\ge0 $, acts on $\ell^2(X,m)$. Moreover, the map $t\mapsto P_t^\omega f$ solves the $\omega$-heat equation for all $f\in\ell^2(X,m)$ and $\omega\in\ell^\infty(X)$.

Following the proof of {\cite[Th.\,2.7]{KellerRose-22a}}, we~obtain the $\ell^2$-mean value inequality {below}. We indicate the necessary changes in the latter article in order to obtain our result as the line of estimates is analogous with minor modifications.
\begin{proposition}[Moser in time and space, \cf {\cite[Th.\,2.7]{KellerRose-22a}}]\label{thm:KR22-subsolution}
Let $ x\in X $, $T\in\RR$,
$n>2$, $\delta\in(0,1]$, $r\geq 128 S$ and constants $\phi,\Phi\geq 1$ be given.
Assume $\rS_{\phi}(n,r/2,r)$ in~$ x$ and the doubling property
\[
m(B_x(r))\leq \Phi\ m(B_x(r/2)).
\]
For all non-negative $\Delta_\omega$-subsolutions $v\geq 0$ on $[T-r^2,T+r^2]\times B_{x}(r)$, we~have
\begin{multline*}
\biggl(\frac{1}{2 \delta (r/2)^2m({B_{x}(r/2)})}\int\limits_{T-\delta (r/2)^2}^{T+\delta (r/2)^2}\sum_{B_{x}(r/2)} mv_t^{2w}\drm t\biggr)^{\sfrac{1}{w}}
\\
\leq
\frac{C_{n,\phi,\Phi}(1+\delta r^2h(\omega))^{\Psfrac{n}{2}+1} }{\delta^{\Psfrac{n}{2}+1} r^{2}m({B_{x}(r)})}
\int\limits_{T- \delta r^2}^{T+\delta r^2}\sum_{B_{x}(r)} mv_t^{2}\drm t,
\end{multline*}
where $C_{n,\phi,\Phi}:=2^{{110}n^2}(\phi\Phi)^{2n}$
and
\[
w=\Bigl(\frac{r}{8S}\Bigr)^{\sfrac{1}{n}}.
\]
\end{proposition}
\begin{proof}
The proof follows along the same lines as the proof of \cite[Th.\,2.7]{KellerRose-22a} with a different choice of radii for the Moser iteration steps. More precisely, using the notation of the latter article, we~choose
\[
\rho_k=\frac{r}{2}(1+2^{-k}), \quad k= 0,\ldots, \eta(r),
\]
with $\eta(r)=\left\lfloor\log_2\Psfrac{r}{16S}\right\rfloor$ iteration steps. To apply the iteration procedure, we~need to check that $\rho_{k}\leq r$ and $ \rho_{k}-\rho_{k+1}>2S $ for $ k= 0,\ldots, \eta(r) $. This is indeed the case since $\rho_{k}\leq r$ is trivial and, for $ k= 0,\ldots, \eta(r)$, we~have, since $S\le r2^{-\eta-4}\leq r 2^{-k-4}$,
\[
\rho_k-\rho_{k+1}-2S={r}{2^{-k-2}}-2S\ge {r}{2^{-k-3}}>0,
\]
which is needed for the iteration procedure, \cf \cite[Prop.\,2.6]{KellerRose-22a}. What remains is to track the constant $C_{n,\phi,\Phi}$. Following the arguments of the proof of \cite[Th.\,2.7]{KellerRose-22a} we estimate
\[
\rho_k^2-\rho_{k+1}^2\geq r^2\,2^{-(k+4)}.
\]
Finally, we~have to replace $2^dC_D$ in the notation of the latter article by $\Phi$ from the doubling property. Since $n>2\geq 1$, we~obtain a constant called $ C_{d,n} $ in \cite[Th.\,2.7]{KellerRose-22a} where in our situation $ d=n $ and $ (1\vee C_{D})=2^{-d}\Phi $
\[
C_{d,n}=\Phi\bigl(\Phi^{\Psfrac{n}{2}+1}\phi^{\sfrac{n}{2}}\bigr)10^{8((n+2)(d+1)+n^2+n)+1}2^{-(1+n/2)d}
\leq [\Phi \phi]^{2n}2^{{{110}n^2}}
\]
and we used $ 2\le n $ as well as $ 10^{3}\leq 2^{10} $.

This leads to a bound on the weighted $\ell^{\tilde p}$-norm of $v$
with $\tilde p(r)=\alpha^{\eta(r)}$, where $\alpha=1+2/n$ and
\[
\eta(r)=\left\lfloor\log_2\frac{r}{16S}\right\rfloor\geq \log_2\frac{r}{16S}-1=\log_2\frac{r}{32S}.
\]
Since the function $f(x)=\ln(1+x)- x\ln 2$ is concave on $ [0,1] $ and \hbox{$ f(0)=f(1)=0 $}, we~infer that $ f \ge 0$ on $ [0,1] $. Hence, employing this for $ x=2/n $, we~obtain $\log_2((n+\nobreak2)/n)\geq 2/n$. Since $r\geq 2^7S$, we~obtain
\[
\alpha^{\eta(r)}\geq \Bigl({\frac{n+2}{n}}\Bigr)^{\log_2\Psfrac{r}{32S}}=\Bigl(\frac{r}{2^5S}\Bigr)^{\log_2\Ppsfrac{n+2}{n}}\geq \Bigl(\frac{r}{2^5S}\Bigr)^{\sfrac{2}{n}} \ge \Bigl(\frac{r}{2^{3}S}\Bigr)^{\sfrac{1}{n}}
= w.
\]
since $ n>2 $.
The claim follows from the monotonicity of weighted $\ell^p$-norms with respect to $p$.
\end{proof}

In order to obtain subsolution estimates, we~use a special case of \cite[Th.\,3.4]{KellerRose-22a} with the choices $\beta=1+2/(n+2)$, $X=\{x\}$ and $\mu=\gamma>0$.

\begin{proposition}[Moser iteration in time, \cf {\cite[Th.\,3.4]{KellerRose-22a}}]\label{thm:constantballsdavies}
Fix $\gamma,\delta>0$, $T\geq 0$, $n>2$, $\beta=1+2/(n+2)$, $k\in\NN_0$, and let $v\geq 0$ be a bounded $\Delta_\omega$-supersolution on the cylinder $[(1-\delta)T,(1+\delta)T]\times \{x\}$.
Then we have
\begin{align*}
\sup_{t\in[(1-\delta/2)T,(1+\delta/2)T]}\!\! v_t(x)^2 \leq G \biggl(\frac{1}{2\delta T}\int\limits_{(1-\delta)T}^{(1+\delta)T}
\gamma\ v_t^{2\beta^k}(x)\ \drm t\biggr)^{\sfrac{1}{\beta^k}}\!\!\!,
\end{align*}
where $G=G_{x,\gamma}(\delta,T,k,n)$ is given by
\[
G= C_{1,n}
\bigl[\left(1+\delta T\Deg(x)\right)
\gamma^{-1}\bigr]^{\sfrac{1}{\beta^{k}}},
\qquad\text{and}\qquad
C_{1,n}=2^{{12} n^2}.
\]
\end{proposition}
\begin{proof}
The proof of \cite[Th.\,3.4]{KellerRose-22a} with the choice $\beta=1+\spfrac2{n+2}=\pspfrac{n+4}{n+2}$ reveals the constant $ C_{\beta} $
\[
C_{\beta}=2^{\bigl(4+\Psfrac{1}{[\ln\pspfrac{n+4}{n+2}]}+\psfrac{n+2}{2}+1\bigr)(n+2)}.
\]
Using the mean value theorem we obtain, since $n\geq 2$,
\[
\frac{1}{\ln[\pspfrac{n+4}{n+2}]}=\frac{1}{\ln(n+4)-\ln(n+2)}\leq \frac{1}{\min_{t\in[n+2,n+4]}1/t}= {\frac{n}{2}+2}\leq {2}n.
\]
Finally, note that
\[
\Bigl(4+\frac{1}{\ln[\pspfrac{n+4}{n+2}]}+\frac{n+2}{2}+1\Bigr)(n+2)
\leq (2n+ {2}n+2n)2n = {12}n^2.\qedhere
\]
\end{proof}

We define for $r\geq0$, $n>2$, $x\in X$, and constants $\phi,\Phi\geq 1$ the error-function $\tilde\Gamma_x(r)
:=
\tilde\Gamma_x(r,n,\phi,\Phi)\geq 0$ by
\[
\tilde\Gamma_x(r)
:=
2^{67n^2}(\phi \Phi(1\vee S))^{2n}
\left(1\vee\Deg(x)
\right)
^{2n\sqrt[n]{{S}/{r}}}.
\]
This error term will later be estimated and split up in the error term $ A\Psi $.

\begin{theorem}\label{thm:l2meanvaluesturm} Let $ x\in X $, $r\geq
128 S$, $n>2$.
Assume that there are constants $\phi,\Phi\geq 1$ such that $\rS_{\phi}(n,r/2,r)$ holds in $ x $ and that the doubling property
\[
m(B_x(r))\leq \Phi\ m(B_x(r/2))
\]
is satisfied. Then for all $\tau\in(0,1]$, $T\geq 0$,
and all non-negative $\Delta_\omega$-solutions $v$ on the cylinder $[T-r^2,T+r^2]\times B_x(r)$ we have
\begin{align*}
v^2 _T(x)
\leq
\frac{\tilde\Gamma_x(r/2)^2(1+\tau r^2h(\omega))^{\Psfrac{n}{2}+1}}
{\tau^{\Psfrac{n}{2}+1}r^{2}m(B_x(r))}
\int\limits_{T-\tau r^2}^{T+ \tau r^2}\sum_{B_x(r)} m\ v_t^{2}\ \drm t.
\end{align*}
\end{theorem}
\begin{proof} Set {$w=\sqrt[n]{r/8S}$} and
\[
\beta=1+\frac{2}{n+2}\quad\text{and}\quad k=\Bigl\lfloor \frac{\ln w}{\ln \beta} \Bigr\rfloor.
\]
Clearly, $w\geq \beta^k{\ge1}$ and $k\geq \ln w/\ln \beta -1$.
Since $L^p$-norms with respect to probability measures are non-decreasing in $p\in[1,\infty]$,
we can use this fact after we applied Proposition~\ref{thm:constantballsdavies} with $\gamma=m(x)/m(B_x(r/2))$ and constant $\delta=\sfrac{\tau(r/2)^2}{T}$ to get
\begin{align*} v_T(x)^2
&\le G\biggl(\frac{1}{2 \tau(r/2)^2m(B_x(r/2))}\int\limits_{T- \tau (r/2)^2}^{T+ \tau (r/2)^2} m(x)\ v_t^{2\beta^{k}}(x)\ \drm t\biggr)^{\sfrac{1}{\beta^{k}}}
\\
&\le G\biggl(\frac{1}{2\tau (r/2)^2m(B_x(r/2))}\int\limits_{T- \tau (r/2)^2}^{T+ \tau (r/2)^2}\sum_{B_x( r/2)} m\ v_t^{2\beta^{k}}\ \drm t\biggr)^{\sfrac{1}{\beta^{k}}}
\\
&\leq G\biggl(\frac{1}{2\tau (r/2)^2m(B_x(r/2))}\int\limits_{T-\tau (r/2)^2}^{T+\tau (r/2)^2}\sum_{B_x( r/2)} m \ v_t^{2w}\drm t\biggr)^{\sfrac{1}{w}}\\
&\leq
C_{n,\phi,\Phi}G\frac{(1+\tau r^2h(\omega))^{\Psfrac{n}{2}+1}}
{\tau^{\Psfrac{n}{2}+1}r^{2}m(B_x(r))}
\int\limits_{T-\tau r^2}^{T+ \tau r^2}\sum_{B_x(r)} m\ v_t^{2}\ \drm t,
\end{align*}
where the last estimate follows by Proposition~\ref{thm:KR22-subsolution} with $\delta=\tau$ which is applicable since $r\ge 128S $. Using the definition of $\beta=1+2/(n+2)$ and the lower bound $k\geq \ln w/\ln \beta -1$, we~get the estimate
\[
\frac{1}{\beta^k}\leq \frac{1}{\beta^{\ln w/\ln \beta-1}}=\frac{\beta}{w}\leq \frac{2}{w}.
\]
We have $C_{n,\phi,\Phi}=2^{{110}n^2}(\phi\Phi)^{2n}$ and {with $ \gamma= m(x)/m(B_x(r/2))$ and $ \delta=\tau r^2 /4T $}
\[
G= {G_{x,\gamma}(\delta,T,k,n)} =2^{{12} n^2} \Bigl[\Bigl(1+\frac{\tau r^2}{4}\Deg(x)\Bigr) \frac{m(B_x(r/2))}{m(x)}\Bigr]^{\sfrac{1}{\beta^k}}.
\]
The estimate $m(B_x(r/2))/m(x)\leq {3}^{n/2}\phi^{n/2}(1\vee (r/2)^n)(1 {\vee}\Deg(x))^{n/2}$ follows from Lemma~\ref{lemma:gammanormalized} by $\rS_\phi(n,r/2,r/2)$. Hence, we~get the statement from
\begin{align*}
C_{n, \phi, \Phi}{G}
&\leq 2^{{122}n^2}(\phi\Phi)^{2n}{3}^{n/w}\phi^{n/w}
\left[
\left(1+{r^2} \Deg(x)
\right)(1\vee r^n)(1 {\vee}\Deg(x))^{\sfrac{n}{2}}
\right]
^{\sfrac{2}{w}}
\\
&\leq 2^{{123}n^2}(\phi\Phi)^{3n}
\left[2(1\vee r^{2n})
(1\vee \Deg(x))^{\psfrac{n+2}{2}}
\right]
^{\sfrac{2}{w}}
\\
&\leq 2^{{124}n^2}(\phi\Phi)^{3n}(1\vee r)^{4n/w}
(
1\vee \Deg(x)
)
^{\sfrac{2n}{w}}
\\
&\leq 2^{{134}n^2}(\phi\Phi)^{3n}
(1\vee S)^{4n}
\left[
1\vee \Deg(x)
\right]
^{{4}n\sqrt[ n]{2S/r}}
\leq {\tilde\Gamma_{x}(r{/2})^2},
\end{align*}
where we used
$ \tau \le 1 $, $w\geq 1$ {and $1/w=\sqrt[n]{8S/r} \leq 2\sqrt[n]{2S/r} $ since $n>2$ as well as}
\[
r^{1/w}=r^{\sqrt[n]{8S/r}}=(8S)^{\sqrt[n]{8S/r}}(r/8S)^{\sqrt[n]{8S/r}}\leq (1\vee 8S)\euler^{n/\euler}\leq (1\vee S) 2^{3+n},
\]
where we used $ x^{\sqrt[n]{1/x}} \le e^{n/e} \le 2^n $ for $ x>0 $, \cf proof of Theorem~\ref{thm:adapteddoubling}, and that we squared $ \tilde \Gamma $ in the last inequality.
\end{proof}

In order to obtain the desired heat kernel bounds from the subsolution estimates for the $\omega$-heat equation, we~will use the following result.

\begin{proposition}[{\cite[Th.\,5.3]{KellerRose-22a}}]\label{theorem:daviesabstractgraph}Let $T>0$,
$Y\subset X$, and
$a, b\colon Y \to [0,\infty)$, $a\leq b$, and $\chi\colon Y\times [0,\infty)\to [0,\infty)$ such that
for all $f\in \ell^2(X)$, $f\geq 0$, $\omega\in\ell^\infty(X)$, and $x\in Y$ we have
\[
\chi(x, h(\omega))^{2}(P_T^\omega f)^2(x)\leq \int_{a(x)}^{b(x)}\Vert P_t^\omega f\Vert_2^2\ \drm t.
\]
Then we have
for all $x,y\in Y$
\begin{multline*}
p_{2T}(x,y)
\leq
\frac
{
(b(x)-a(x))^{\sfrac{1}{2}}(b(y)-a(y))^{\sfrac{1}{2}}
\exp\left(\frac{b(x)+b(y)-2T}{2}\nu(\rho_{xy},2T)\right)
}
{\chi\big(x,\nu(\rho_{xy},2T)\big)\chi\big(y,\nu(\rho_{xy},2T)\big)}
\\
\quad \cdot\exp\big(-\Lambda(a(x)+a(y))-(2T/S^2)\zeta(\rho_{xy}S/2T)\big),
\end{multline*}
where
\[
\rho_{xy}:=\rho(x,y),\quad \nu(r,t):=2 S^{-2}\biggl(\sqrt{1+\frac{r^2S^2}{t^2}}-1\biggr).
\]
\end{proposition}

The next result is a variant of \cite[Th.\,6.1]{KellerRose-22a} for varying dimensions and an error term depending only on the degree and reciprocal measure of {one vertex} rather than means of these quantities in a growing ball.

\begin{theorem}\label{thm:main1sturm}
Let $x,y\in X$, $\phi\geq 1$ be a constant, and
$ R_2\geq 4R_1\geq 4096S$. For $z\in\{x,y\}$, let $ n_z\colon[R_1,R_2]\to(2,\infty) $ be given and set for $\tau\in [4R_1,R_2]$
\[
N=N_z(\tau)=\sup_{[\tau/4,\tau]}n_z.
\]
If $\rS_\phi(n,R_1,R_2)$ holds in $\{x,y\}$,
then $\rG_{A\Psi}(N,4R_1,R_2)$ holds in $\{x,y\}$,
where
\[
\Psi_z(N,\tau)=\left(1\vee \Deg(z)
\right)^{7N^2\sqrt[N]{S/\tau}}\quad\text{and}\quad
A(N)=2^{54N^3}\phi
^{3N^2}(1\vee S)^{11N^2}.
\]
\end{theorem}
\begin{proof}
The proof is divided into two parts which we explain before we get into the details. First we use that the Sobolev inequality implies volume doubling, Theorem~\ref{thm:adapteddoubling}. This is then used in Theorem~\ref{thm:l2meanvaluesturm} to conclude a mean value inequality. In~turn we use this together with Proposition~\ref{theorem:daviesabstractgraph} to show after an appropriate choice of the involved constants an estimate of the form
\begin{multline*}
p_{t}(x,y)
\leq 2^{4{\tilde N_{xy}}}
\tilde \Gamma_x(r/2)\tilde \Gamma_y(r/2)
\\
\cdot
\frac{
\left(1\vee {S^{-2}}{\left(\sqrt{t^2+\rho_{xy}^2S^2}-t\right)}\right)^{\sfrac{\tilde N_{xy}}{2}}}
{\sqrt{m(B_{x}(r))m(B_{y}(r))}}\,
\euler^{-\Lambda (t- r^2)-(t/S^2)\zeta(\rho_{xy}S/t)},
\end{multline*}
for $ r=\sqrt{t/2}\wedge R_{2} $, where $\tilde \Gamma $ are the error terms in Theorem~\ref{thm:l2meanvaluesturm}, $ \tilde N_{xy}$ is a dimension function, $S$ is the global jump size and $ \rho_{xy}=\rho(x,y) $ is the intrinsic metric. From there the second part is then rather technical to further estimate the involved terms to their desired final form.

For the first part, let $z\in\{x,y\}$, $ T\ge R_{1}^{2} $ and fix $ r=\sqrt{T}\wedge R_2 \in [R_{1},R_{2}] $.
To ease notation we denote
\[
\tilde N:=\tilde N_{z}(r):={\sup_{[r/2,r]}n_z.}
\]
First from Lemma~\ref{lem:sobdim}, we~get that $\rS_\phi (n,R_1,R_2)$ in $z$, implies $\rS_{\phi} (\tilde N, r/2,r)$ in $z$ for every $ r\in[2R_1,R_2] $.
Then, since $r\in[S,R_2/2]$, Theorem~\ref{thm:adapteddoubling} yields the doubling property with constant
\[
C_z(r)=2^{9 \tilde N^2}\phi^{\sfrac{\tilde N}{2}}(1\vee S)^{4\tilde N}\left(1\vee\Deg(x)\right)^{2\tilde N\sqrt[\tilde N]{S/r}}.
\]
The Sobolev inequality $\rS_{{\phi}}(N,r/2,r)$ and volume doubling property are the assumptions of Theorem~\ref{thm:l2meanvaluesturm}, where the doubling constant is denoted by $\Phi $, \ie we have $ \Phi=C_{z}(r) $. We apply Theorem~\ref{thm:l2meanvaluesturm} to the function $ v $ given by
\[
(t,x)\mapsto P^\omega_t f(x)
\]
for $f\in \ell^2(X)$, $f\geq 0$, $\omega\in\ell^\infty(X)$, which is an $\omega$-solution on $[0,\infty)\times X$. Hence, from Theorem~\ref{thm:l2meanvaluesturm}, we~obtain
for all $\delta\in(0,1]$, $T\geq \delta r^2$, $z\in\{x,y\}$
\begin{align*}
P_T^\omega f(z)^2
\leq
\frac{\tilde\Gamma_{z}(r/2)^2(1+\delta r^2h(\omega))^{\Psfrac{\tilde N}{2}+1}}{\delta^{\Psfrac{\tilde N}2+1}r^2m(B_{z}(r))}
\int\limits_{T-\delta r^2}^{T+\delta r^2}\Vert\Eins_{B_{z}(r)} P_t^\omega f\Vert_2^{2}\drm t,
\end{align*}
where $
\tilde\Gamma_x(r)
:=
2^{67\tilde N^2}(\phi C_{x}(r)(1\vee S))^{2\tilde N}
\left(1\vee\Deg(x)
\right)
^{2\tilde N\sqrt[\tilde N]{{S}/{r}}}. $
To finally apply
Proposition~\ref{theorem:daviesabstractgraph} we choose the parameters.
We have $T-\delta r^2\geq 0$. We set
\[
a(z)=T-\delta r^2,\quad b(z)=T+\delta r^2,\quad
r(z)=r,
\]
and $\chi(z, h(\omega))$ via
\[
\chi(z,h(\omega))^{-2}=\frac{\tilde\Gamma_{z} (r/2)^2}{\delta^{\Psfrac{\tilde N}2+1}r^{2}
m(B_{z}(r))}(1+\delta r^2h(\omega))^{\Psfrac{\tilde N}{2}+1}.
\]
Proposition~\ref{theorem:daviesabstractgraph} yields
for $T\geq 4R_1^2$, and $ t=2T $
the estimate
\begin{align*}
p_{t}(x,y)&=p_{2T}(x,y)
\\&
\leq \frac
{
(b(x)-a(x))^{\sfrac{1}{2}}(b(y)-a(y))^{\sfrac{1}{2}}
\exp\left(\frac{b(x)+b(y)-2T}{2}\nu(\rho_{xy},2T)\right)
}
{\chi\left(x,\nu(\rho_{xy},2T)\right)\chi\left(y,\nu(\rho_{xy},2T)\right)}
\\
&\hspace*{4cm} \cdot\exp\left(-\Lambda(a(x)+a(y))-(2T/S^2)\zeta(\rho_{xy}S/2T)\right)
\\
&
=
\frac{2
\tilde\Gamma_{x} (r/2)\tilde\Gamma_{y}(r/2)\delta^{-\sfrac{\tilde N_{xy}}{2}}}
{\sqrt{m(B_{x}(r))m(B_{y}(r))}} \left(1+\delta r^2\nu(\rho_{xy},t)\right)^{\Psfrac{\tilde N_{xy}}{2}+1}
\\
&\hspace*{2.5cm}\cdot
\exp\left(\delta r^2\nu(\rho_{xy},t)\right)
\exp\bigl(-\Lambda (t-2\delta r^2)-(t/S^2)\zeta(\rho_{xy}S/t)\bigr),
\end{align*}
where $\tilde N_{xy}=(\tilde N_x(r)+\tilde N_y(r))/2$, $\nu(r,t)=2 S^{-2}\left(\sqrt{1+{r^2S^2}{t^{-2}}}-1\right) $, and $\rho_{xy}=\rho(x,y)$. Choosing
\[
\delta=\frac{1}{2}\wedge \frac{1}{t\nu(\rho_{xy},t)},
\]
we obtain, since $r=\sqrt{t/2}\wedge R_2\leq \sqrt t$ and $ \exp(1)\leq 4 \leq 2^{\tilde N} $,
\begin{multline*}
p_{t}(x,y)
\leq 2^{4{\tilde N_{xy}}}
\tilde \Gamma_x(r/2)\tilde \Gamma_y(r/2)
\\
\cdot \frac{
\left(1\vee {S^{-2}}{\left(\sqrt{t^2+\rho_{xy}^2S^2}-t\right)}\right)^{\sfrac{\tilde N_{xy}}{2}}}
{\sqrt{m(B_{x}(r))m(B_{y}(r))}}\,
\euler^{-\Lambda (t-t\wedge R_2^2)-(t/S^2)\zeta(\rho_{xy}S/t)}.
\end{multline*}

The second part, which is the rest of the proof, is devoted to estimate the appearing terms which is done in three steps: We estimate the volume terms, the dimension term and the remaining error terms.
We abbreviate $\tau=\sqrt t\wedge R_2$. Recall $ r=\sqrt{t/2}\wedge R_2 $ and let $ z\in \{x,y\} $.

In the first step, we~estimate the volume terms $1/ m(B_{z}(r))\leq C_{z}(\tau)/m(B_{z}(\tau)) $ via volume doubling:
Since $ t\geq 16R_1^2$ and $ R_{2}\ge 4R_{1} $, we~clearly have the trivial inequality $r= \sqrt{t/2} \wedge R_2\geq \tau/2\geq 2R_1$. The doubling property leads to
\begin{align*}
\frac{1}{m(B_{z}(r))}
=
\frac{1}{m(B_{z}(\sqrt {t/2}\wedge R_2))}
\leq
\frac{1}{m(B_{z}(\tau/2))}
\leq
\frac{C_z(\tau)}{m(B_{z}(\tau))}
\end{align*}
with
{$ C_z(r)=2^{9 \tilde N^2}\phi^{\sfrac{\tilde N}{2}}(1\vee S)^{4\tilde N}\left(1\vee\Deg(x)\right)^{2\tilde N\sqrt[\tilde N]{S/r}} $.}

For the second step, we~turn to the dimension terms. Note $\tau/2\leq r\leq \tau$. We have
\begin{align*}
\max_{[r,\tau]}\tilde N_{z}&= \max_{R\in [r,\tau]}\Vert n_z\Vert_{[R/2,R] }\leq \Vert n_z\Vert_{[\tau/4,\tau] }= N_{z}(\tau)=:N.
\end{align*}
This yields the estimate $ \tilde N_{xy}=(\tilde N_{x}(\tau)+\tilde N_{y}(\tau))/2 \leq( N_{x}(\tau)+ N_{y}(\tau))/2=N_{xy} $ needed for the correction term
\begin{align*}
\left(1\vee {S^{-2}}{\left(\sqrt{t^2+\rho_{xy}^2S^2}-t\right)}\right)^{\sfrac{\tilde N_{xy}}{2}}\leq \Bigl(1\vee {S^{-2}}{\left(\sqrt{t^2+\rho_{xy}^2S^2}-t\right)}\Bigr)^{\sfrac{N_{xy}}{2}}.
\end{align*}

We are left to estimate $ 2^{{\tilde N} }\tilde \Gamma_{z}(r/2) \sqrt{C_{z}(\tau)}\leq A^{N}_{z}(\tau)\Psi_{z}(\tau) $ for each $ z\in \{x,y\} $ which is the third and final step.
We start with the term
\[
C_z(r)=2^{9 \tilde N^2}\phi^{\sfrac{\tilde N}{2}}(1\vee S)^{4\tilde N}\left(1\vee\Deg(x)\right)^{2\tilde N\sqrt[\tilde N]{S/r}}
\]
arising from volume doubling. This term can be estimated by
\begin{align*}
\max_{[r,\tau]}C_{z} \leq {2^{9 N^2}\phi^{\sfrac{N}{2}}(1\vee S)^{4N}\left(1\vee\Deg(x)\right)^{2N\sqrt[N]{S/r}}},
\end{align*}
where we used the short hand $ N=N_{z}(\tau) $, from above. We
recall
\[
\tilde\Gamma_z(r)
:=
2^{67\tilde N^2}(\phi C_{z}(r)(1\vee S))^{2\tilde N}
\left(1\vee\Deg(z)
\right)
^{2\tilde N\sqrt[\tilde N]{{S}/{r}}}
\]
and $\tau/2\leq r\leq \tau$ to estimate
\begin{align*}
\tilde \Gamma_{z}&(r/2)\sqrt{C_{z}(\tau)}
=2^{67\tilde N^2}\big(\phi\cdot C_{z}(r)(1\vee S)\big)^{2\tilde N}
\left(1\vee\Deg(z)\right)^{2\tilde N\sqrt[\tilde N]{2{S}/{r}}} \sqrt{C_{z}(\tau)}
\\
&\leq 2^{67 N^2}\Bigl(2^{9 N^2} \phi^{\Psfrac{N}{2}+1}(1\vee S)^{5N}(1\vee\Deg(z))^{2N\sqrt[N]{S/\tau}} \Bigr)^{2 N}
\cdot\left(1\vee\Deg(z)\right)^{4N\sqrt[ N]{{S}/\tau}}
\\ &\hspace*{2cm} \cdot 2^{4 N^2}\phi^{\sfrac N4}(1\vee S)^{2N}\left(1\vee\Deg(z)\right)^{N\sqrt[N]{S/\tau}}
\\
&=
2^{18N^3+143 N^2/2}\phi^{N^{2}+\sfrac{9N}{4}} (1\vee S)^{10N^2+2N} \left(1\vee \Deg(z)\right)^{(4N^2+5N)\sqrt[n]{S/r}}\\
& \leq 2^{54 N^{3}}\phi^{3N^{2}}(1\vee S)^{11N^2}\left(1\vee \Deg(z)
\right)^{7N^2\sqrt[N]{S/\tau}},
\end{align*}
where we used $ N\ge 2 $. Hence, setting $ \Psi_z(\tau)=\left(1\vee \Deg(z)
\right)^{7N^2\sqrt[N]{S/\tau}}$
and $ A(N)=2^{54N^3}\phi
^{3N^2}(1\vee {S})^{11N^2} $
reveals the estimate {$ 2^{2{\tilde N}}\tilde \Gamma_{z}(r/2)\sqrt{C_{z}(\tau)} \leq A(N_{z}(\tau))\Psi_{z}(\tau)$} which finishes the proof.
\end{proof}

\section{On-diagonal bounds imply Sobolev inequalities}\label{section:heatsob}

{In this section} we show that the existence of a family of operators satisfying certain contractivity properties provides weak Sobolev inequalities. Then we use these weak Sobolev inequalities to derive Sobolev inequalities. Finally, we~apply these results to obtain Sobolev inequalities from upper heat kernel bounds. The results and proofs are inspired by the corresponding results on manifolds \cite{SaloffC-01,Varopoulos-85}. Special attention is paid to the influence of the local geometry of the underlying graph on the behavior of the involved constants.

Often we will use the convention
\[
\Vert f\Vert_W:={\|f1_{W}\|_{\infty}=}\sup_W\vert f\vert
\]
for a function $f\colon W\to \RR$ defined on some subset $ W $ of $ X $ or $ X\times [0,\infty]$.
\begin{lemma}[Weak Sobolev]\label{lem:weakS}Let $C_1,C_2,n>0$, $0\leq r_1\leq r_2$,
and $(Q_{r})_{r\in[r_1, r_2]}$ be a family of operators with domain including $\cC_c(X)$ such that for $f\in\cC( B(r_2))$ and all $r\in [ r_1, r_2]$
\[
\Vert Q_{r}f\Vert_\infty\leq C_1 r^{-n}\Vert f\Vert_1, \quad \Vert f-Q_{r}f\Vert_2\leq C_2 r\Vert |\nabla f|\Vert_2.
\]
Then,
\[
\sup_{\lambda>0}\lambda^{2(1+1/n)}m(f>\lambda)
\leq 12 C_2^2
\Bigl(C_1\vee \Bigl(r_1^n \Bigl\Vert\frac1m\Bigr\Vert_{B(r_2)}\Bigr)\Bigr)^{{2}/{n}}
\Bigl( \Vert|\nabla f|\Vert_2^2+\frac{1}{r_2^2}\Vert f\Vert_2^2\Bigr)\Vert f\Vert_1^{2/n}.
\]
\end{lemma}

\begin{proof}If $f=0$ in $B(r_2)$ there is nothing to prove, so assume $f\neq 0$ in $B(r_2)$. If~$\lambda>\Vert f\Vert_\infty$, then $\{f>\lambda\}=\emptyset$ and there is nothing to prove, so assume $\lambda \leq \Vert f\Vert_\infty$. We always have
\[
\lambda^2 m(f>\lambda)=\sum_X m \lambda^2\Eins_{\{f>\lambda\}}\leq \Vert f\Vert_2^2.
\]
Let $\mu>0$ to be chosen later. We distinguish between two cases for $\lambda\geq 0$.

First, if $\lambda\leq \mu C_1 r_2^{-n}\Vert f\Vert_1$, we~get
\begin{align*}
\lambda^{2(1+1/n)}m(f>\lambda)=\lambda^2m(f>\lambda)\lambda^{2/n}&\leq (\mu C_1)^{2/n} r_2^{-2}\Vert f\Vert_2^2\Vert f\Vert_1^{2/n}
\\
&\leq (\mu C_1)^{2/n}\left( \Vert|\nabla f|\Vert_2^2+r_2^{-2}\Vert f\Vert_2^2\right)\Vert f\Vert_1^{2/n}.
\end{align*}
Second, let $\lambda> \mu C_1 r_2^{-n}\Vert f\Vert_1$, and set
\[
r_f:=(\mu\ C_1\Vert f\Vert_1\lambda^{-1})^{1/n}.
\]
First, for the upper bound on $r_f$, note that the current lower bound of $\lambda$ yields
\[
r_f=
\Bigl(\frac{\mu C_1\Vert f\Vert_1}{\lambda}\Bigr)^{1/n}
\leq
\Bigl(\frac{\mu C_1\Vert f\Vert_1}{\mu C_1 r_2^{-n}\Vert f\Vert_1}\Bigr)^{1/n}=r_2.
\]
Now, we~estimate $r_f$ from below.
Since
$\lambda\leq\Vert f\Vert_\infty$,
we obtain
\begin{align*}
r_f
=
\Bigl(\frac{\mu C_1\Vert f\Vert_1}{\lambda}\Bigr)^{1/n}
&\geq
\Bigl(\frac{\mu C_1\Vert f\Vert_1}{\Vert f\Vert_\infty}\Bigr)^{1/n}.
\end{align*}
If we choose
\[
\mu
\geq \frac{r_1^n}{C_1} \frac{\Vert f\Vert_\infty}{\Vert f\Vert_1},
\]
then
$r_f\geq r_1$.
Hence, we~can apply our assumptions to $r_f\in[ r_1,r_2]$.

For later use, we~need to choose $\mu$ such that $\vert Q_{r_f} f\vert <\lambda/2$ on $B(r_2)$: This is satisfied if we assume $\mu\geq 3$. Indeed, assuming the existence of $x\in B(r_2)$ with $\vert Q_{r_f} f\vert(x) \geq \lambda/2$, the definition of $r_f$ and our assumption on $\Vert Q_{r}f\Vert_\infty$ yield
\[
\frac{3}{2}C_1\Vert f\Vert_1 r_f^{-n}
\leq \frac{\mu}{2}C_1\Vert f\Vert_1 r_f^{-n}=\frac{\lambda}{2}\leq\vert Q_{r_f}f\vert(x)\leq \Vert Q_{r_f}f\Vert_\infty\leq C_1{r_f}^{-n}\Vert f\Vert_1,
\]
a contradiction.

The restrictions on $\mu>0$ above lead us to the choice
\[
\mu= 3\Bigl(1\vee \Bigl(\frac{r_1^n}{C_1} \frac{\Vert f\Vert_\infty}{\Vert f\Vert_1}\Bigr)\Bigr).
\]
The triangle inequality yields
\[
\{f> \lambda \}\subset \{\vert f-Q_{r}f\vert \geq\lambda/2\}\cup\{\vert Q_{r}f\vert\geq \lambda/2\}
\]
for all $r\geq 0$ and $\lambda>0$. Hence, since $\vert Q_{r_f}f\vert <\lambda/2$ on $B (R_2)$, we~obtain using our second assumption and the definition of $ r_{f} =\left(\sfrac{\mu\ C_1\Vert f\Vert_1}{\lambda}\right)^{1/n}$ and \hbox{$ \mu=3\bigl(1\vee \bigl(\frac{r_1^n}{C_1} \frac{\Vert f\Vert_\infty}{\Vert f\Vert_1}\bigr)\bigr) $}
\begin{align*}
m(f>\lambda)
\leq m\left(\vert f-Q_{r_f}f\vert \geq \lambda/2\right)
&\leq \Bigl(\frac{2}{\lambda}\Bigr)^2\Vert f-Q_{r_f}f\Vert_2^2
\leq
\Bigl(\frac{2}{\lambda}\Bigr)^2C_2^2r_f^2\Vert| \nabla f|\Vert_2^2
\\
&=
\Bigl(\frac{2}{\lambda}\Bigr)^2
C_2^2
\left(\left(
\mu C_1\Vert f\Vert_1\lambda^{-1}\right)^{1/n}\right)^2\Vert| \nabla f|\Vert_2^2
\\
&= \frac{4\cdot 3^{\sfrac{2}{n}} C_2^2}{\lambda^{2(1+1/n)}}
\Vert |\nabla f|\Vert_2^2\big[
(C_1\Vert f\Vert_1)\vee \left(r_1^n \Vert f\Vert_\infty
\right)\big]^{\sfrac{2}{n}}.
\end{align*}
In order to obtain the desired right-hand side, we~use
\[
\Vert f\Vert_\infty=\max_{B(r_2)}|f|
\leq \sup_{B(r_2)}\frac1m \Vert f\Vert_1=\Bigl\Vert\frac1m\Bigr\Vert_{B(r_2)} \Vert f\Vert_1. \hfill\qedhere
\]
\end{proof}

In order to derive the Sobolev inequality from the lemma above, we~provide a version of the proof of \cite[Th.\,4.1]{BakryCLS-95} and \cite{Barlow-book,Delmotte-97}.

\begin{theorem}[Sobolev]\label{thm:SobM}
Let $C_1,C_2>0$, $n>2$, $0\leq r_1\leq r_2$,
and $(Q_r)_{r\in[r_1,r_2]}$ a family of operators defined on $\cC_c(X)$ such that for all $f\in\cC(X)$ with $\supp f\subset B(r_2)$ and all $r\in [r_1,r_2]$
\[
\Vert Q_rf\Vert_\infty\leq C_1 r^{-n}\Vert f\Vert_1, \quad \Vert f-Q_rf\Vert_2\leq C_2 r\Vert| \nabla f|\Vert_2.
\]
Then, we~have for all $f\in\cC(X)$ with $\supp f\subset B(R_2)$
\begin{align*}
\Vert f\Vert_{\spfrac{2n}{n-2}}^2
\leq 2^{8+\spfrac{2n}{n-2}}C_2^2
\Bigl(C_1\vee \Bigl(r_1^n\Bigl\Vert\frac1m\Bigr\Vert_{B(r_2)} \Bigr) \Bigr)^{\sfrac{2}{n}}
\Bigl(
\Vert |\nabla f|\Vert_2^2+\frac{1}{r_2^2}\Vert f\Vert_2^2 \Bigr).
\end{align*}
\end{theorem}
\begin{proof}
We choose the following partition of identity for $ k\in \mathbb{Z} $
\[
\Phi_k\colon [0,\infty)\to[0,\infty),\quad t\longmapsto (t-2^k)_+\wedge 2^k
=
\begin{cases}
0,& t<2^k,\\
t-2^k, &2^k\leq t<2^{k+1},\\
2^k,& t\geq 2^{k+1}.
\end{cases}
\]
Indeed, for $t\geq 0$, we~can calculate by choosing $l\in \ZZ$ such that $2^l\leq t<2^{l+1}$
\[
\sum_{k\in \mathbb{Z}}\Phi_k(t)=t-2^l+\sum_{k=-\infty}^{l-1}2^k=t.
\]
If $t,s\geq0$, using $\sum_k a_k^2\leq (\sum_k \vert a_k\vert)^2$, this yields (w.l.o.g.~$t\geq s$, \ie $ \phi_{k}(t)\ge \phi_{k}(s) $)
\[
(t-s)^2= \Bigl(\sum_{k\in \mathbb{Z}}(\Phi_k(t)-\Phi_k(s)) \Bigr)^2
\geq \sum_{k\in \mathbb{Z}}(\Phi_k(t)-\Phi_k(s))^2.
\]
If $f\ge 0$, then we set $f_k:=\Phi_k\circ f$, $k\in \ZZ$. Hence, by the estimate above
\[
\sum_{k\in \mathbb{Z}}\Vert \vert \nabla f_k\vert\Vert_2^2\leq \Vert \vert \nabla f\vert\Vert_2^2.
\]
Moreover,
since $0\leq f_k\leq 2^k$ and $(\Phi_k)_k$ is a partition of the identity, we~have
\begin{align*}
\sum_{k\in \mathbb{Z}} \Vert f_k\Vert_2^2
\leq \Vert f\Vert_2^2.
\end{align*}
Hence, if we let
\[
W(f):=\Vert| \nabla f|\Vert_2^2+\frac{1}{r_2^2}\Vert f\Vert_2^2,
\]
then
\[
W(f_k)
\leq W(f), \quad k\in\ZZ.
\]
Abbreviate
\[
C_M:=2^{4}C_2^2
\Bigl(C_1\vee \Bigl(r_1^n \Bigl\Vert\frac1m \Bigr\Vert_{B(r_2)} \Bigr) \Bigr)^{\sfrac{2}{n}},
\]
and
\[
N(f):=\sup_{k\in\ZZ} 2^k m(f\geq 2^k)^{\sfrac{1}{q}}, \qquad q=\frac{2n}{n-2}.
\]
\begin{claim}We have {for all $ f\ge 0 $}
\[
N(f)^{2}
\leq
2^{q}C_M
W(f).
\]
\end{claim}
We postpone the claim to the end of the proof and show how to derive the desired from the claim. Hence, we~use
$ \{f_k\geq 2^k\}=\{f\geq 2^{k+1}\} $, the inequality $\sum \vert a_k\vert^{q/2}\leq \left(\sum \vert a_k\vert\right)^{q/2}$ which is applicable since $q/2\geq 1$ and the claim to estimate
\begin{multline*}
\Vert f\Vert_q^q
=\sum_{X}m\sum_{k\in\mathbb{Z}}\vert f\vert^q\Eins_{\{2^k\leq f< 2^{k+1}\}}
\leq \sum_{k\in\mathbb{Z}}2^{q(k+1)}\sum_{X}m \Eins_{\{2^k\leq f< 2^{k+1}\}}
\\
\leq \sum_{k\in\mathbb{Z}}2^{q(k+1)}m(f\geq 2^k)
=\sum_{k\in\mathbb{Z}}2^{q(k+2)}m(f\geq 2^{k+1})
=2^{2q}\sum_{k\in\mathbb{Z}}2^{qk}m(f_k\geq 2^{k})
\\
{\le}2^{2q}\sum_{k\in\mathbb{Z}}N(f_k)^q
\leq 2^{2q}\Bigl(\sum_{k\in\mathbb{Z}}N(f_k)^2\Bigr)^{q/2}
\leq 2^{2q}2^{q^2/2}C_M^{q/2}\Bigl(\sum_{k\in\mathbb{Z}}
W(f_k)\Bigr)^{q/2},
\end{multline*}
\ie
\[
\Vert f\Vert_q^2
\leq 2^{4+q}C_M\sum_{k\in\mathbb{Z}}W(f_k)\leq 2^{4+q}C_M W(f).
\]
For general $f$, decompose $f=f_+-f_-$ and obtain the above inequalities for $f_+$ and~$f_-$. Since $ f_{+} $ and $ f_{-} $ are orthogonal in $\ell^2(X,m)$ and $ \||\nabla f_{+}|\|^{2} +\||\nabla f_{-}|\|^{2}\leq \||\nabla {f}|\|^{2} $, the claim follows for $f$.

\noindent
\emph{Proof of the claim.} Denote $\tau=1+1/n$.
Since $f_k\leq 2^k$ and $ \supp f_{k} \subset \{f\geq 2^k\} $, we~have {with $ q=2n/(n-2) $}
\[
\Vert f_k\Vert_1\leq 2^km(f\geq 2^k)
=2^{k\left(1-q\right)}(2^{k}m(f\geq 2^k)^{\sfrac{1}{q}})^{q}
\leq
2^{k\left(1-q\right)}N(f)^{q}.
\]
{The weak Sobolev inequality,} Lemma~\ref{lem:weakS}, applied to $\lambda=2^k$ and $f_k$ yields together with {$ W(f_{k})\leq W(f) $ and} the above estimate
\[
2^{k2\tau}m(f_k\geq 2^k)
\leq C_MW(f_k)\Vert f_k\Vert_1^{\sfrac2n}
\leq
C_M2^{k\left(1-q\right)\sfrac2n}W(f)N(f)^{q\Psfrac2n}.
\]
We have
$ \{f\geq 2^{k+1}\}=\{f_k\geq 2^k\} $ and $ (k+1)q-k2\tau+k\left(1-q\right)\Psfrac2n=q$, and hence,
\begin{align*}
\bigl(2^{k+1}m(f\geq 2^{k+1})^{\sfrac{1}{q}}\bigr)^{q}
&=2^{(k+1)q-k2\tau}2^{k2\tau} m(f_k\geq 2^k)
\\
&\leq
C_M2^{(k+1)q-k2\tau}2^{k\left(1-q\right)\Psfrac2n}W(f)N(f)^{q\Psfrac2n}
\\
&=
C_M2^{q
}W(f)N(f)^{q\Psfrac2n}.
\end{align*}
The definition of $N(f)$ yields
\[
N(f)^{q}
\leq
C_M2^{q}
W(f)N(f)^{q\Psfrac2n}.
\]
{Dividing by $N(f)^{2q/n}$}, using $q-q\Psfrac2n=2$ yields the claim {and finishes the proof}.
\end{proof}

In the following we obtain the Sobolev inequality from heat kernel bounds in balls depending on the local geometry of the graph. The original idea of proof goes back to \cite{Varopoulos-85}. The argument is nowadays standard. We include the argument to track the constants and for the sake of completeness.

\begin{theorem}\label{lem:heattoS}
Let $0\leq r_1\leq r_2$, and assume for all $x,y\in B(r_2)$
\[
p_{r^2}(x,y)\leq Cr^{-n}, \quad r\in [r_1,r_2].
\]
Then, for all {$f\in\cC( B(r_{2}))$}, we~have
\begin{align*}
\Vert f\Vert_{\spfrac{2n}{n-2}}^2
\leq 2^{8+\spfrac{2n}{n-2}}
\Bigl(C\vee \Bigl(r_1^{n} \Bigl\Vert\frac1m\Bigr\Vert_{B(r_2)}\Bigr)\Bigr)^{\sfrac{2}{n}}
\Bigl(\Vert \vert\nabla f\vert\Vert_2^2+\frac{1}{r_2^2}\Vert f\Vert_2^2\Bigr).
\end{align*}
\end{theorem}
\begin{proof}
We need to check the assumptions of Theorem~\ref{thm:SobM} for $Q_r=P_{r^{2}}=e^{-r^2\Delta}$.
Clearly, we~have with $ B=B(r_2) $
\[
\Vert P_{r^2}\Vert_{\ell^1(B)\to\ell^\infty(B)}=\sup_{x,y\in B}p_{r^2}(x,y)\leq C r^{-n}, \quad r\in [r_1,r_2],
\]
what is the first assumption of Theorem~\ref{thm:SobM}.
Now, we~check the second assumption. Fix $f\in \cC(X)$ with $\supp f\subset B$ and let $ t=r^{2} $. Then by self-adjointness of $P_t$
we have
\begin{align*}
\Vert f-P_tf\Vert_2^2=\Vert f\Vert_2^2+\Vert P_tf\Vert_2^2-2\langle P_tf,f\rangle
=\Vert f\Vert_2^2+\Vert P_tf\Vert_2^2-2\Vert P_{t/2}f\Vert_2^2.
\end{align*}
By the contraction property of the heat semigroup, \ie $\Vert P_tf\Vert_2^2\leq \Vert P_{t/2}f\Vert_2^2$, and the fundamental theorem of calculus we obtain
\begin{align*}
\Vert f-P_tf\Vert_2^2
\leq \Vert f\Vert_2^2-\Vert P_{t/2}f\Vert_2^2=-\int_0^{t/2}\frac{\drm }{\drm s}\Vert P_sf\Vert_2^2\drm s
=-2\int_0^{t/2}\Bigl\langle \frac{\drm }{\drm s}P_sf,P_sf\Bigr\rangle\drm s.
\end{align*}
Since $s\mapsto P_sf$ solves the heat equation $ -\frac{\drm }{\drm s}P_sf=\Delta P_sf $, Green's formula yields
\begin{align*}
\Vert f-P_tf\Vert_2^2
\leq
2\int_0^{t/2}\langle \Delta P_sf,P_sf\rangle\drm s
=2\int_0^{t/2}\Vert |\nabla P_sf |\Vert_2^2\drm s,
\end{align*}
Analogously, the inequality
\[
\Vert |\nabla P_sf|\Vert_2^2\leq \Vert|\nabla f|\Vert_2^2
\]
can be seen by using the fundamental theorem, the heat equation, and Green's formula to obtain
\[
\Vert |\nabla P_sf|\Vert^2- \Vert|\nabla f|\Vert^2
=-\int_0^s \Vert \Delta P_sf\Vert_2^2\drm s
\leq 0.
\]
Thus, we~get since $r^{2}=t$ and $ Q_{r}=P_{r^{2}} $
\begin{align*}
\Vert f-Q_{r}f\Vert_2^2=\Vert f-P_{t}f\Vert_2^2
\leq
2\int_0^{t/2}\Vert |\nabla P_sf|\Vert_2^2\drm s
\leq t\Vert |\nabla f|\Vert_2^2=r^2\Vert |\nabla f|\Vert_2^2.
\end{align*}
The assumptions of {Theorem~\ref{thm:SobM}} are satisfied for $r\in[r_1,r_2]$ and the claim follows with $C_1=C$ and $C_2=1$.
\end{proof}

\section{{Gaussian} upper bounds and {volume} doubling imply Sobolev}\label{section:GS}

In this section we show how to derive a Sobolev inequality from Gaussian upper bounds \eqref{G} and volume doubling \eqref{V}. We pursue two lines of argument. The first more classical approach, Theorem~\ref{thm:heattosobolevgeneral}, going back to \cite{Varopoulos-85} collects all the remnants of the estimates within the Sobolev constant. This is later used for the normalizing measure since in this case these remnants can be uniformly bounded. However, in the general unbounded case this is not feasible. In~this case we mitigate the unbounded remains by choosing a variable dimension function to get a general preliminary Sobolev inequality Theorem~\ref{thm:heattosobolevgeneralimproving}. In~this theorem, we~still have a free parameter called $ \gamma $ which is then appropriately chosen in the next section.

Both approaches use a sequence of technical lemmas given below, which are an elaboration on the general strategy given in \cite{SaloffC-01}. First we show how to bound the measure of a ball in terms of the measure of any other ball via chaining and the doubling condition. This is the only place in the paper where we use that $ \rho $ is a path metric.
\begin{lemma}[Comparing balls]\label{lem:ballcomparison}
Let {$r\ge 8S $}, $ o\in X $, $ d>0$, $\Psi\geq 1$ be constants,
and assume $\rV_\Psi (d,r/4,r)$ in $B_{o}(r)$. Then, for all $x,y\in B_o(r)$, we~have
\[
m(B_x(r))
\leq
2^{18d}\,
\Psi^{9}
\ m(B_{y}(r)).
\]
\end{lemma}
\begin{proof}
If $x, y\in B(r)$ such that $B_x(r/4)\cap B_y(r/4)\neq \emptyset$, then clearly $B_x(r/4)\subset B_y(r)$. By $\rV_\Psi(d,r/4,r)$, we~have
\[
m(B_x(r))
\\
\leq 4^d\Psi\ m(B_x(r/4))\leq 4^d\Psi\ m(B_y(r)).
\]

Let $x,y\in B(r)$ and
denote by $p$ a path $x=z_0\sim z_1\sim \ldots \sim z_l= y$ realizing $\rho(x,y)$ which exists due to local finiteness \cf \cite[Ch.\,12.2]{KellerLW-21}. We construct a sequence of vertices $ (z_{k_i}) $ on $ p $ inductively and denote $ B_{i}=B_{z_{k_{i}}}(r/4) $. Set $k_0:=0$. For given $ k_{i} $ and $ B_{i} $, we~choose the smallest index $ k =1,\ldots,l$ such that $ z_{k+1} $ is not in $\bigcup_{j=0}^{i} B_{j} $ but $ z_{k} $ is in $\bigcup_{j=0}^{i} B_{j} $ and set $ k_{i+1}={k} $. If~there is no such $ z_{k+1} $ on the path we choose $ k_{i+1}=l $ and we set $ L=i+1 $. {Observe that $ z_{k}=z_{k_{i+1}}\in B_{i}=B_{z_{k_{i}}}(r/4) $ since if $ z_{k}\in B_{j} $ for $ j\neq i $ and using that the path $ z_0\sim\ldots\sim z_{l} $ realizes $ \rho(x,y) $, we~get a contradiction
\begin{align*}
\frac{r}{4}\ge \rho(z_{k},z_{k_{j}})= \rho(z_{k},z_{k_{i}})+ \rho(z_{k_{i}},z_{k_{j}})>\frac{r}{4}.
\end{align*}
}Moreover, since $ z_{k+1}\not\in B_{i}= B_{z_{k_{i}}}(r/4) $ and {$ \rho(z_{k+1},z_{k_{i+1}})=\rho(z_{k},z_{k+1})\leq S $, we~have
\begin{align*}
\rho ( z_{k_i}, z_{k_{i+1}})\ge \rho(z_{k_{i}},z_{k+1})-\rho(z_{k+1},z_{k_{i+1}})>\frac{r}{4}-S > 0
\end{align*}}
for $ i=0,\ldots,L-1 $. Thus, we~have
\begin{align*}
r\ge \rho(x,y)=\sum_{i=0}^{L-1}\rho(z_{k_{i}},z_{k_{i+1}}) \ge (L-1)\Bigl(\frac{r}{4}-S\Bigr).
\end{align*}
Since $ r\ge 8S $, we~obtain
\begin{align*}
L\le \frac{5r-4S}{r-4S}= 5 + \frac{16S}{r-4S}\leq 9.
\end{align*}
At the same time, we~have
$ B_{i}\cap B_{i+1} \neq \emptyset$ since {$ z_{k_{i+1}}\in B_{i}$ as shown above $z_{k_{i+1}}\in B_{i+1}=B_{z_{k_{i+1}}}(r/4) $} for $ i=0,\ldots, L $.
Iterating the estimate in the beginning of the proof along $ (z_{k_{i}}) $ we obtain
\begin{align*}
m(B_{x}(r))\leq 4^{Ld}\Psi^{L}m(B_{y}(r))\leq 4^{9d}\Psi^{9}m(B_{y}(r))
\end{align*}
which yields the claim.
\end{proof}

Next, we~need the notion of on-diagonal bounds which plays only a technical role in our considerations.

\begin{definition} Let $B\subset X$, $ r_{2}\ge r_{1}\ge 0 $, and $\Psi\colon B\times [r_1,r_2]\to (0,\infty)$. The \emph{on-diagonal estimate} $\rO_{\Psi}(r_1,r_2)$ is satisfied in $B$ if we have
\[
p_{\rho^2}(x,x)\leq \frac{\Psi(x,\rho)^2}{m(B_x(\rho))},\quad x\in B,\; \rho\in[r_1,r_2].
\]
\end{definition}
Gaussian upper bounds $\rG_\Psi(n,r_1,r_2)$ in $B$ obviously imply $\rO_\Psi(r_1,r_2)$ in $B$ for all $[r_1,r_2]$. Indeed, the dimension $ n $ does not appear in the on-diagonal bounds and while $\rG_\Psi(n,r_1,r_2) $ is an assumption on all $ t=r^{2}\ge r_{1} ^{2}$, the on-diagonal bounds are only required for $ r \in[r_{1},r_{2}]$ (which is why no $ \sqrt{t}\wedge r_{2}=r\wedge r_{2} $ terms appear).

\begin{lemma}[Uniform on-diagonal bounds]\label{lem:heattosob}
Let
$0\leq {4}r_1\leq r$, $8S\leq r$, and constants $\Psi\geq 1$, $d>2$ such that $\rO_\Psi(r_1,r)$
and $\rV_\Psi (d,r_1,r)$ hold in $B(r)$.
Then for $x,y\in B(r)$, $\sigma\in[r_1,r]$, we~have
\[
p_{\sigma^2}(x,y)
\leq
2^{18d}\,\Psi^{20}
\frac{r^{d}}{m(B(r))}\,\sigma^{-{d}}.
\]
\end{lemma}
\begin{remark}The proof of Lemma~\ref{lem:heattosob} requires the comparison of volumes of balls with same radius but different centers, Lemma~\ref{lem:ballcomparison}, which needs that the intrinsic metric is a path metric. Imposing $\rV_\Psi (d,r_1,2r)$ in $B(r)$ in Lemma~\ref{lem:heattosob} instead allows to drop this restriction on the metric.
\end{remark}
\begin{proof} Since $8S\leq r$,
and $\rV_\Psi (d,r_1,r)$ hold in $B(r)=B_o(r)$, we~use Lemma~\ref{lem:ballcomparison} to infer
\[
m(B(r))\leq 2^{18d}\,\Psi^9\ m(B_x(r)), \quad x\in B(r).
\]
For any $\sigma\in[r_1,r]$ and $x\in B(r)$, we~obtain from $\rO_{\Psi}(r_1,r)$ and $\rV_\Psi (d, r_1,r)$ in $x$
\begin{align*}
p_{\sigma^{2}}(x,x)\leq \frac{\Psi^2}{m(B_x(\sigma))}
\leq
\frac{\Psi^3 }{m(B_x(r))}\Bigl(\frac{r}{\sigma}\Bigr)^{d}
\leq
\frac{2^{18d}\Psi^{20}}{m(B(r))} \Bigl(\frac{r}{\sigma}\Bigr)^{{d}}.
\end{align*}
For $\sigma\in[r_1,r]$ and $x,y\in B(r)$, we~infer
\[
p_{\sigma^2}(x,y)\leq \sqrt{p_{\sigma^2}(x,x)p_{\sigma^2}(y,y)}
\leq
\frac{2^{18d}\Psi^{20}}{m(B(r))}\Bigl(\frac{r}{\sigma}\Bigr)^{{d}},
\]
where the first inequality follows directly from the semigroup identity and Cauchy-Schwarz inequality.
Thus, $\sigma\leq r$ implies the claim for all $n\geq d$.
\end{proof}

\begin{lemma}[On-diagonal bounds and volume doubling imply Sobolev]\label{lem:heattosob3}
Let
$0\leq {4} r_1\leq r$, $8S\leq r$ and constants $\Psi\geq 1$, $d>2$, and assume $\rO_\Psi(r_1,r)$ and
$\rV_\Psi (d,r_1,r)$ hold in $B(r)$.
Then for all $n\geq d$ we have $\rS_{\tilde \phi}(n,r,r)$ in $o$, where $\tilde \phi =\tilde \phi(r_1,r) $ is given by
\[
\tilde \phi(r_1,r)=2^{44+\spfrac{2n}{n-2}}
\Bigl[\Psi^{20}\vee r_1^n\frac{m(B(r))}{r^{n}}\Bigl\Vert\frac1m\Bigr\Vert_{B(r)}\Bigr]^{\sfrac{2}{n}}.
\]
\end{lemma}

\begin{proof}
Lemma~\ref{lem:heattosob} yields for all $n\geq d$, $\sigma\in[r_1,r]$ and $x,y\in B(r)$
\[
p_{\sigma^2}(x,y)
\leq
C
\sigma^{-{n}}, \quad \text{where}\quad C=2^{18n}\,\Psi^{20}\frac{r^{n}}{m(B(r))}.
\]
Thus, if $n\geq d$,
Theorem~\ref{lem:heattoS} yields
for all $f\in\cC_{c}( B(r))$
\begin{align*}
\Vert f\Vert_{\spfrac{2{n}}{{n}-2}}^2
&
\leq 2^{8+\spfrac{2n}{n-2}}\Bigl(C\vee \Bigl(r_1^{n} \Bigl\Vert\frac1m\Bigr\Vert_{B(r)}\Bigr)\Bigr)^{{\sfrac{2}{n}}}\Bigl(
\Vert \vert\nabla f\vert\Vert_2^2+\frac{1}{r^2}\Vert f\Vert_2^2\Bigr),
\end{align*}
which yields the result since $2^{8+\psfrac{2n}{n-2}}(C\vee (r_1^{n} \left\Vert\tfrac1m\right\Vert_{B(r)}))^{{\sfrac{2}{n}}}\leq \tilde\phi(r_{1},{r}) $.
\end{proof}

From now, we~let the error terms be given by functions rather than constants. For $0\leq r_1\leq r$, let
\[
\Psi\colon X\times(2,\infty)\times[r_1,r]\to[1,\infty),\qquad (x,n,\tau)\longmapsto \Psi_{x}(n,\tau),
\]
be given. For $ x\in X $ and $ r,s\ge0 $, we~denote the cylinder
\[
Q(r,s)=B_{x}(s)\times [r,s].
\]
Without adapting the dimensions we end up with the following.

\begin{theorem}[Sobolev -- fixed dimension]\label{thm:heattosobolevgeneral}Assume $ n\!>\!2 $ is a constant,
\hbox{$0\!\leq\! {4} r_1\!\leq\! r_2$}, $8S\leq r_2$, $x\in X$ and that $ \rO_\Psi(r_1,r_2)$ and $\rV_\Psi (n,r_1,r_2)$ hold in $B(r_2)=B_{x}(r_{2})$.
Then we have $\rS_\phi(n, {4}r_1,r_2)$ in $x$, where
\[
\phi(r)=2^{44+\spfrac{2n}{n-2}}\Bigl[{\Vert \Psi\Vert_{Q(r_1,r)}^{20}}\vee \Bigl( r_1^{n}\frac{m(B(r))}{r^{n}}\Bigl\Vert\frac1m\Bigr\Vert_{B(r)}\Bigr) \Bigr]^{\sfrac{2}{n}}.
\]
\end{theorem}

\begin{proof}
For $r\in[r_1,r_2]$, Lemma~\ref{lem:heattosob3} implies $\rS_{\tilde \phi}(n,r,r)$ in $x$ with $ \tilde\phi =\tilde{\phi}(r_{1},r)$.
\end{proof}

From Theorem~\ref{thm:heattosobolevgeneral} we immediately get \eqref{G} \& \eqref{V} $ \Rightarrow $ \eqref{S} with the respective Sobolev constant above. This is acceptable if the norms of $ \Psi $ and $ 1/m $ stay bounded. For graphs with unbounded geometry this cannot be expected beyond the normalizing measure.

However, as one can see, a~large dimension $ n $ potentially mitigates large norms of $ \Psi $ and $ 1/m $ as it enters as a large root in the Sobolev constant. We pursue this strategy in the following. First, in Lemma~\ref{lem:variabledim1}, we~estimate the dimension such that the Sobolev constant $ \phi $ stays uniformly bounded. Secondly, in Theorem~\ref{thm:heattosobolevgeneralimproving}, we~choose radii $ r_{1} $ in dependence of $ r_{2} $ such that the Sobolev dimension $ n(r_{2}) $ converges as $ r_{2}\to \infty $.

We now also allow for a variable dimension
\[
d\colon X\times [r_1,r]\to (2,\infty)
\]
in the volume doubling property.

\begin{lemma}[Choosing a dimension]\label{lem:variabledim1}Let
$0\leq {4}r'\leq r$, $8S\leq r$, $ x\in X $, let $\Psi\geq 1$, $d>2$ be constants, and assume $\rO_\Psi(r',r)$ and
$\rV_\Psi (d,r',r)$ hold in $B(r)=B_{x}(r)$.
Then, for all $\gamma\geq 1$, we~have $\rS_{\hat \phi}(n,r,r)$ in $o$, where
\[
\hat{\phi}=\hat{\phi}(d)=2^{47+\spfrac{2d}{d-2}}\,
\gamma^{\sfrac{2}{d}}
\]
and the dimension $ n $ can be chosen such that
\[
n\geq d\vee
\ln
\Bigl(1\vee
\frac{\Psi^{20}}{\gamma}\vee
\Bigl[\frac{m(B(r))}{\gamma} \Bigl\|\frac1m \Bigr\|_{B(r)}
\Bigr]^{\sfrac{1}{\ln ({r}/{r'})}}
\Bigr).
\]
\end{lemma}

\begin{proof}
Since $0\leq {4}r'\leq r$, Lemma~\ref{lem:heattosob3}
yields for all $n\geq d$ the Sobolev inequality $\rS_{\tilde{\phi}}(n,r,r)$ in $o$ with
\[
\tilde{\phi}=\tilde{\phi}(r',r)=2^{44+\spfrac{2n}{n-2}}
\Bigl[\frac{\Psi^{20}}\gamma\vee (r')^n\frac{m(B(r))}{\gamma r^{n}} \Bigl\Vert\frac1m \Bigr\Vert_{B(r)} \Bigr]^{\sfrac{2}{n}}\gamma^{\sfrac{2}{n}},
\]
where we snuck in $\gamma\geq 1$. The second entry of the maximum can be estimated by~$1$,~\ie
\[
(r')^n\,\frac{m(B(r))}{\gamma r^n} \Bigl\Vert\frac1m \Bigr\Vert_{B(r)}\leq 1
\]
if and only if
\[
\ln \Bigl(\frac{m(B(r))}\gamma\Bigl\Vert\frac1m\Bigr\Vert_{B(r)}\Bigr)^{\sfrac{1}{{\ln ({r}/{r'})}}}\leq n.
\]
Hence, under this condition,
\[
\tilde\phi(r',r)\leq 2^{44+\spfrac{2n}{n-2}}
\Bigl[1\vee \frac{\Psi^{20}}\gamma\Bigr]^{\sfrac{2}{n}}\gamma^{\sfrac{2}{n}}.
\]
Note that if $t\geq 1$, then we have for all $n\geq \ln t$ the estimate $t^{1/n}\leq t^{1/\ln t}=\euler$. Hence, if we further require
\[
n\geq \ln\Bigl(\frac{\Psi^{20}}\gamma\vee 1\Bigr),
\]
and using $e^2\leq 2^3$, we~obtain
\[
\tilde{\phi}(r',r)\leq 2^{44+\spfrac{2n}{n-2}}\,
\euler^2\,\gamma^{\sfrac{2}{n}}\leq 2^{47+\spfrac{2n}{n-2}}\,
\gamma^{\sfrac{2}{n}}\leq 2^{47+\spfrac{2d}{d-2}}\,
\gamma^{\sfrac{2}{d}}=\hat{\phi}(d)
\]
since the function $ n\mapsto 2n/(n-2) $ is decreasing and $ n\ge d $.
Putting the conditions on~$n$ together yields the claim.
\end{proof}

Now we are in the position to choose specific radii. In~contrast to the lemma above, we~assume now that $\Psi\colon X\times [R_1,R_2]\to [1,\infty)$ is a function. The parameter $\gamma$ is used to adjust the error terms later.

\begin{theorem}[Sobolev -- variable dimension]\label{thm:heattosobolevgeneralimproving}
Let $32S\leq {4} R_1\leq R_2$, $ o\in X $, $n\colon X\times [R_1,R_2]\to (2,\infty)$, and assume that $\rG_{\Psi}(n,R_1,R_2)$ and $\rV_{\Psi} (n,R_1,R_2)$ hold in $B(R_2)=B_{o}(R_{2})$.
Then for all functions $\gamma\colon [4R_1,R_2]\to [1,\infty)$ we have $\rS_{\phi}(n',{4} R_1,R_2)$ in $o$, where
\begin{align*}
\phi(r)=2^{49+\spfrac{2 N(r)}{N(r)-2}}\,\gamma(r)^{\sfrac{2}{N(r)}}
\end{align*}
with $N(r)=\|n(r)\|_{B(r)}$ and the dimension $ n' $ can be chosen
\begin{align*}
n'(r)\geq N(r)\vee \ln\biggl[1\vee
\frac{{\Vert\Psi\Vert_{Q(r',r)}^{20}}}{\gamma(r)}\vee \Bigl(\frac{m(B(r))\Vert \frac1m\Vert_{B(r)}}{\gamma(r)}\Bigr)^{\sfrac{1}{\ln(r/r')}}
\biggr]
\end{align*}
with
\[
r'= R_1\vee\Bigl( \frac{r}{4}\wedge \frac{1}{4}(\ln r)^{2N(r)}\Bigr).
\]
\end{theorem}

\begin{proof}Gaussian bounds $\rG_\Psi(n,R_1,R_2)$ in $B(R_2)$ yield on-diagonal bounds~$\rO_\Psi(r'\mkern-4mu,r)$ in $B(r)$ if $R_1\leq {r}'\leq r\leq R_2$. Similarly, $\rV_\Psi (n,R_1,R_2)$ in $B(R_2)$ yields $\rV_\Psi (N,r',r)$ in $B(r)$ if $R_1\leq r'\leq r\leq R_2$. With these observations the
result follows immediately from Lemma~\ref{lem:variabledim1} applied to the interval $[r',r]$ if $R_1\leq {r}'$ and ${4} r'\leq r$ for all $r\in[R_1,R_2]$ which is trivially satisfied.
\end{proof}

The theorem above shows how to deduce \eqref{S} from \eqref{G} and \eqref{V} with a parameter function $\gamma$ and general $\Psi$. Next, we~present a corollary where we choose a particular~$\gamma$, given the function $\Psi$, as
\[
\Psi_z(n,\tau)=\left(1\vee \Deg(z)
\right)^{7n^2\sqrt[n]{S/\tau}}
\]
for some $n>2$ as in Theorem~\ref{thm:main1sturm}. This result is more general than Theorem~\ref{thm:generaldeg1}\eqref{thm:generaldeg1ii}, since it also includes \eqref{S} for small radii.
\begin{corollary}\label{cor:sobolevall}
Let
\[
32S\leq {4} R_1\leq R_2,\quad o\in X,\quad A\colon X\times [R_1,R_2]\to[1,\infty),\quad {n}\colon X\times [R_1,R_2]\to(2,\infty)
\]
and assume that $\rG_{A\Psi}(n,R_1,R_2)$, $\rV_{\!A\Psi} (n,R_1,R_2)$, $\rL_A(n,R_1,R_2)$ hold in $B(R_2)=B_{o}(R_{2})$.
Then we have $\rS_{\phi'}(n',{4} R_1,R_2)$ in $o$, where
\begin{align*}
\phi'(r)&= 2^{85+\spfrac{2 N}{N-2}}\tilde A^{20},
\\
n_o'(r)&=
N \Bigl(1+\frac{\ln(r'/(1\wedge r))}{\ln(r/r')} +{28
0 (1\vee S) N{\frac{\ln (1\vee r)}{\sqrt[ N]{r'}}}}\Bigr)\Bigl(1+ \frac{1}{2}\frac{\ln\Vert 1\vee \Deg\Vert_{B(r)}} {\ln(1\vee r)}\Bigl),
\\
N&= \Vert n\Vert_{Q(r',r)},\qquad \tilde A= \Vert A\Vert_{Q(r',r)},\qquad \tilde N(r)=\|n\|_{Q(R_1\vee (\ln r)/4,r)},
\\
r'&= R_1\vee\Bigl( \frac{r}{4}\wedge \frac{1}{4}(\ln r)^{2\tilde N(r)}\Bigr).
\end{align*}
\end{corollary}
\begin{proof}
For any $\gamma\colon[R_1,R_2]\to[1,\infty)$, Theorem~\ref{thm:heattosobolevgeneralimproving} (together with Lemma~\ref{lem:sobdim}) yields $\rS_{\tilde \phi}(\tilde n, 4R_1,R_2)$ in $o$ for $ \tilde \phi\ge \phi^{\gamma} $ and $ \tilde n\ge n^{\gamma} $ where
\begin{gather*}
\phi^{\gamma}(r)=2^{49+\spfrac{2 N}{N-2}}\gamma(r)^{\sfrac{2}{N}},
\quad
n^{\gamma}(r)= N\vee \ln\Bigl[1\vee
\frac{T_1}{\gamma(r)}\vee \Bigl(\frac{T_2}{\gamma(r)}\Bigr)^{\sfrac{1}{\ln(r/r')}}
\Bigr],
\\
T_1:={\Vert A\Psi\Vert_{Q(r',r)}^{20}},\qquad T_2:=m(B_o(r))\Bigl\|\frac1m\Bigr\|_{B_o(r)}.
\end{gather*}
In order to bound $n^{\gamma}$ from above, we~need to estimate $T_1$ and $T_2$ from above and choose $\gamma$ appropriately. We start with $T_1$ and abbreviate $
D :=\Vert (1\vee\Deg(x))\Vert_{B(r)}$. Then,
\begin{align*}
\Vert \Psi\Vert_{Q(r',r)}
=\sup_{(x,\tau)\in Q(r',r)} (1\vee\Deg(x))^{7n_z(\tau)^2\sqrt[n_z(\tau)]{S/\tau}}
\leq
D^{7 N^2\sqrt[ N]{S/r'}}.
\end{align*}
We conclude with $ \tilde A= \Vert A\Vert_{Q(r',r)}$
\[
T_1
\leq
\Vert A\Vert_{Q( r',r)}^{20}\Vert \Psi\Vert_{Q(r',r)}^{20}
\leq\tilde A^{20} D^{140 N^2\sqrt[ N]{S/r'}}.
\]
Now, we~estimate $T_2$. First, since distance balls are finite by assumption, there exists $x\in B_o(r)$ such that
\[
\Bigl\Vert \frac1m\Bigr\Vert_{B_o(r)}=m(x)^{-1}.
\]
{As $\rV_{\!A{\Psi}}(n,R_1,R_2)$ implies $\rV_{\!\tilde A{\tilde\Psi}}( N,R_1,R_2)$} with {$ \tilde \Psi=\|\Psi\|_{Q(r',r)} $}, Lemma~\ref{lem:ballcomparison} gives
\[
m(B_o(r))
\leq
2^{18 N}
\tilde A^9{\tilde\Psi^{9}}
\ m(B_{x}(r)).
\]
We infer from $\rL_{{A}}(n,R_1,R_2)$ the estimate
\[
\frac{m(B_{x}(r))}{m(x)}
\leq {A(1\vee r)^{n_x(r)}(1\vee\Deg(x))^{\sfrac{n_x(r)}{2}}}
\leq {\tilde A}{((1\vee r^2)
D)^{\sfrac{N}{2}}}.
\]
Hence, using $ D \geq 1$, and
{$\tilde \Psi=\Vert \Psi\Vert_{Q(r',r)}
\leq
D^{7 N^2\sqrt[ N]{S/r'}}$}, we~get
\begin{align*}
T_2
=\frac{m(B_{o}(r))}{m(x)}
\leq 2^{18 N}
\tilde A^9\tilde\Psi^9
\frac{m(B_{x}(r))}{m(x)}
\leq 2^{18 N}\tilde A^{10}
((1\vee r^2) D)^{\Psfrac{N}{2}+63 N^2\sqrt[ N]{S/r'}}.
\end{align*}
Choose
\[
\gamma(r)= {2^{18 N}\tilde A^{{20}}}
\]
to obtain, with $ \tilde A\ge 1 $,
\[
\frac{T_1}{\gamma(r)}
\leq
\frac{\tilde A^{20} D^{140 N^2\sqrt[ N]{S/r'}}}{\gamma(r)}
\leq D^{140 N^2\sqrt[ N]{S/r'}}
\]
and
\[
\frac{T_2}{\gamma(r)}
\leq
\frac{2^{18 N}\tilde A^9((1\vee r^2) D)^{\Psfrac{N}{2}+{63 N^2\sqrt[ N]{S/r'}}}}{\gamma(r)}
\leq ((1\vee r^2) D)^{\Psfrac{N}{2}+{63 N^2\sqrt[ N]{S/r'}}}.
\]
Hence, we~obtain since $ D\geq 1$ and with $ \iota=1/\ln(r/r') \le 1/\ln {4}\le {1}$ as $ r/r' \ge {4}$
\begin{align*}
n^\gamma(r)&=
N\vee \ln\Bigl[1\vee
\frac{T_1}{\gamma(r)}\vee \Bigl(\frac{T_2}{\gamma(r)}\Bigr)^{\iota}
\Bigr]
\\
&\leq
N\vee \ln \left[ D^{140 N^2\sqrt[ N]{S/r'}}\vee \left(((1\vee r^2) D)^{\Psfrac{N}{2}\iota+{63\iota N^2\sqrt[ N]{S/r'}}}\right)\right]
\\
&\leq
N\vee \ln \Bigl[ ((1\vee r^2) D)^{\Psfrac{N}{2}\iota+{140 N^2(1\vee S)\sqrt[ N]{1/r'}}}\Bigr]
\\
&=
N\Bigl[1\vee \ln (1\vee r)\Bigl({{\iota}+{280 N(1\vee S)\sqrt[ N]{\sfrac {1}{r'}}}}\Bigr) \Bigl( 1+ \frac{1}{2}\frac{\ln ( D)}{\ln (1\vee r)}\Bigr)
\Bigr]=n_o'(r),
\end{align*}
where the last equality follows since
\begin{multline*} \iota \ln (1\vee r)= \frac{\ln (1\vee r)}{\ln (r/r')} =1 + \frac{\ln (1\vee r) -\ln (r/r') }{\ln (r/r')}
\\
=1 + \frac{\ln(r'(1\vee r)/r)}{\ln(r/r')} =1 + \frac{\ln(r'/(1\wedge r))}{\ln(r/r')}.
\end{multline*}
Finally, we~{use} the function
\[
\phi^{\gamma}(r) =2^{49+\spfrac{2 N}{N-2}}\gamma(r)^{\sfrac{2}{N}}=2^{85+\spfrac{2 N}{N-2}} A^{40/ N}\leq \phi'(r)
\]
for the $ \gamma(r)= 2^{18 N}\tilde A^{{20}}$ chosen above.
This yields $\rS_{\phi'}( n', 4R_1,R_2)$ in $o$ since $ \phi'\ge \phi^{\gamma} $ and finishes the proof.
\end{proof}

\section{Proof of the main theorems}\label{section:remaining}

We are now ready to prove the main theorems. {We deal with the case of the normalizing measure and the case of general measure including the counting measure separately.
The reason for this distinction is that the proof of Theorem~\ref{thm:normalizedglobal} follows much closer the classical approach and is therefore significantly simpler. Furthermore, we~present also a localized version of Theorem~\ref{thm:normalizedglobal} in Theorem~\ref{thm:equivnormalized}.
}
Afterward, we~will deal with the general case for which we first show a result where the bounds depend on vertex degrees and reciprocals of measures, Theorem~\ref{thm:general}. To reduce the bounds to depend on the vertex degree alone, we~incorporate the local regularity property \eqref{L} to show Theorem~\ref{thm:generaldeg1}. Finally, Theorem~\ref{thm:countingglobal} is an immediate corollary of Theorem~\ref{thm:generaldeg1}.

\subsection*{The normalizing measure}
First, we~will present our results about the normalizing measure $m=\deg$. In~this special case the Laplacian $ \Delta $ is always a bounded operator on $ \ell^{2}(X,\deg) $. Furthermore, the combinatorial distance is then an intrinsic metric which we will always use in this case. In~particular, the jump size is constantly equal to $ 1 $ for this metric.
\begin{theorem}[Normalizing measure, localized version]\label{thm:equivnormalized}Let $ m=\deg $, {$n>2$}, $ R_2\geq 8 R_1\geq 512$.
\begin{enumeratei}
\item If there is a constant $\phi\geq 1$ such that $\rS_\phi (n,R_1,R_2)$ holds in $B\subset X$, then there {is a constant} $\Psi=\Psi(\phi,n,R_1)\geq 1$ such that $\rG_{\Psi}(n,4R_1,R_2)$ and $\rV_{{\Psi}}(n,R_2)$ hold in~$B$.
\item
If there {is a constant $\Psi\geq 1$} such that $\rG_{\Psi}(n,R_1,R_2)$ and $\rV_{{\Psi}}(n,R_2)$ hold in $B_o(R_2)$, then there is a constant $\phi=\phi(n,R_1,\Psi)\geq 1$ such that $\rS_{\phi}(n,4R_1,R_2)$ holds in $o$.
\end{enumeratei}
\end{theorem}

\skpt
\begin{proof}[Proof of Theorem~\ref{thm:equivnormalized}]
\begin{enumeratei}
\item
We have to show that $\rS_{\phi}(n,R_1,R_2)$ in $B$ implies properties $\rV_{{\Psi}}(n,R_2)$ and $\rG_{\Psi}(n,4R_1,R_2)$ in $B$ for appropriate $\Psi>0$.
Corollary~\ref{cor:adapteddoubling} yields $\rV_{{\Psi}}(n,R_2)$ with ${{\Psi}}=2^{{9}n^2}\phi^{n^2} R_1^n$.

From Theorem~\ref{thm:main1sturm} we infer $\rG_\Psi(n,4R_1,R_2)$ with
\[
\Psi_z(n,\tau)=2^{53n^3}\phi
^{3n^2}
\]
as the jump size satisfies $ S=1 $ for the normalizing measure and the combinatorial graph distance.

\item
We have to show that conditions $\rV_{{{\Psi}}}(n,R_2)$ and $\rG_{\Psi}(n,R_1,R_2)$ imply property $\rS_{\phi}(n,4R_1,R_2)$ for appropriate $\phi>0$. By Theorem~\ref{thm:heattosobolevgeneral} we have $\rS_{\phi}(n,4R_1,R_2)$ in $o$ with
\[
\phi(r)=2^{44+\spfrac{2n}{n-2}}\Bigl({\Psi^{20}}\vee R_1^{n}\frac{m(B_o(r))}{r^{n}}\Bigl\|\frac{1}{m}\Bigr\|_{B_o(r)}\Bigr)^{\sfrac{2}{n}}.
\]
Since distance balls are finite and $m(x)=\deg(x)$, there exists $x\in B_o(r)$ such that
\[
\Bigl\|\frac{1}{m}\Bigr\|_{B_o(r)}=\frac{1}{m(x)}=\frac{1}{\deg(x)}=\frac{1}{m(B_x(1/2))}.
\]
Lemma~\ref{lem:ballcomparison} yields $m(B_o(r))\leq 2^{18n}{{\Psi^9}}m(B_x(r))$. Together with the volume doubling property $\rV_{{\Psi}}(n,1/2,R_2)$, we~obtain
\begin{align*}
R_1^{n}\,\frac{m(B_{o}(r))}{r^{n}}\,\Bigl\|\frac{1}{m}\Bigr\|_{B_o(r)}
=\frac{R_1^{n}}{\deg(x)}\frac{m(B_{o}(r))}{r^{n}}
\leq \frac{2^{18n}{{\Psi^9}} R_1^{n}}{m(B_x(1/2))}\frac{m(B_{x}(r))}{r^{n}}
&\leq 2^{19n}{{\Psi^{10}}}R_1^n.
\end{align*}
Hence, $\phi$ is bounded above and the claim follows.\qedhere
\end{enumeratei}
\end{proof}

\subsection*{{General measures}}
In the case of general measures, correction terms in \eqref{G}, \eqref{V}, and \eqref{S} also depend on the vertex degree, which may be unbounded.
{Hence, in this case, we~cannot hope for uniform bounds as for the normalizing measure.
} In order to include unbounded vertex degree into our results, we~employ the results from the preceding sections. They lead to the following general version of the equivalence of Sobolev inequalities and heat kernel bounds involving varying dimensions.

Let $R_{2}\ge {4}R_1\geq0$. For a dimension function $n\colon X\times [R_1,R_2]\to (2,\infty)$ and $ x\in X $ and $ r\in [R_{1},R_2] $, we~let the supremum over annuli be given as
\begin{align*}
N_x(r)=\Vert n_x\Vert_{[r/4,r]}.
\end{align*}
Furthermore, the volume doubling and the Gaussian correction term are given in terms of
\[
\Psi_x(n,r)
=\left(1\vee\Deg(x)\right)^{7n\sqrt[n]{S/(r\vee R_1)}}
\cdot \begin{cases}
((1\vee r^2)(1\vee \Deg(x)))^{\sfrac{n}{2}},& r<R_1,
\\
1,& \text{else.}
\end{cases}
\]
For a constant $ \phi\ge 1$, let
\[
A_x'(r)
=(2\phi)^{\sfrac{n_x(R_1)^2}{2}}\, 2^{54N_x(r)^3}\phi^{3N_x(r)^2}(1\vee S)^{11N_x(r)^2}.
\]
For $o\in X$ and a dimension function $ n $, we~consider the supremum over the space-time cylinder
$ Q(r',r)=B_o(r)\times [r',r] $ and let
\begin{align*}
N(r)=\Vert n\Vert_{Q(r',r)},
\end{align*}
\[
r'= R_1\vee\Bigl( \frac{r}{4}\wedge \frac{1}{4}(\ln r)^{2\tilde N(r)}\Bigr)\qquad {\tilde N(r)=2\|n\|_{Q(R_1\vee (\ln r)/4,r)}}.
\]
Next, for
$\gamma\colon X\times [R_1,R_2]\to (0,\infty)$, we~define a function $ \phi^{\gamma} $, which will play the role of the Sobolev constant, by
\begin{equation*}
\phi^{\gamma}(r)=2^{{49}+\spfrac{2N(r)}{N(r)-2}}\gamma(r)^{\sfrac{2}{N(r)}}.
\end{equation*}
Finally, we~choose a dimension function $ n^{\gamma}$ which satisfies the following inequality
\begin{align*}
n^{\gamma}(r)\geq N(r)\vee \ln\biggl[1\vee
\frac{{\Vert A'\Psi\Vert_{Q(r',r)}^{20}}}{\gamma(r)}\vee \Bigl(\frac{m(B(r))\Vert \tfrac1m\Vert_{B(r)}}{\gamma(r)}\Bigr)^{\sfrac{1}{\ln(r/r')}}
\biggr].
\end{align*}

With this notation the main result in its most general form is now just a consequence of Theorems~\ref{thm:fulldoubling}, \ref{thm:main1sturm}, and \ref{thm:heattosobolevgeneralimproving}.

\begin{theorem}[Most general case]\label{thm:general}Let
$ R_2\geq {4}R_1\geq 0$, {$o\in X$} $\gamma\colon X\times [R_1,R_2]\to (0,\infty)$, and $ \phi\ge 1$ be a constant.
\begin{enumeratei}
\item If $\rS_\phi (n,R_1,R_2)$ holds in $B\subset X$, $R_1\geq 1024S$, then
$B$ satisfies $\rV_{\!A'{\Psi}}(N,R_2)$ and $\rG_{A'\Psi}(N,4R_1,R_2)$.
\item
Assume $\rV_{\!A'{\Psi}} (N,R_2)$ and $\rG_{A'\Psi}(N,R_1,R_2)$ hold in $B_o(R_2)$, $R_1\geq 8S$. Then the property $\rS_{\phi^{\gamma}}(n^{\gamma},4 R_1,R_2)$ holds in $o$.
\end{enumeratei}
\end{theorem}

We are now in the position to prove the all main theorems presented in the introduction.
\begin{proof}[Proof of Theorem~\ref{thm:normalizedglobal}]
This follows from Theorem~\ref{thm:equivnormalized} with $R_2=\infty$.
\end{proof}

To obtain Theorems~\ref{thm:countingglobal} and \ref{thm:generaldeg1} presented in the introduction, we~will incorporate the local regularity condition \eqref{L} into Theorem~\ref{thm:general}. The resulting correction functions and dimensions appearing in \eqref{G}, \eqref{V}, and \eqref{S} depend only on the vertex degree.
We will first prove Theorem~\ref{thm:generaldeg1} before we reduce it to the counting measure case Theorem~\ref{thm:countingglobal}.

\skpt
\begin{proof}[Proof of Theorem~\ref{thm:generaldeg1}]
\par
\eqref{thm:generaldeg1i}
Theorems~\ref{thm:adapteddoubling} and \ref{thm:main1sturm} yield properties $\rV_{\!A\Psi}(n,R_1,R_2)$ and $\rG_{A \Psi}(N,4R_1,R_2)$ in~$B$, with
\[
\Psi_x(N,r)
=\left(1\vee\Deg(x)\right)^{7N\sqrt[N]{S/r}}
\]
and
\[
A(N)
=2^{54N^3}\phi^{3N^2}(1\vee S)^{11N^2}.
\]
\eqref{thm:generaldeg1ii}
This follows from Corollary~\ref{cor:sobolevall}.
\end{proof}

\begin{proof}[Proof of Theorem~\ref{thm:countingglobal}] Since $ \Deg=\deg/m =\deg$ for $ m=1 $ and $S=1$, this follows from {Theorem~\ref{thm:generaldeg1} with $R_2=\infty$ choosing $ n $ and $ \phi $ to be constant}. To obtain the specific form of $ n_o' $ given in Theorem~\ref{thm:countingglobal}, we~estimate $ n_o' $ from Theorem~\ref{thm:generaldeg1} from above. First of all, pick
$r_1\ge 4r_0 \vee 1$ such that $(\ln r)^{4n} \le r$ for all $r\ge r_1$. This is possible since $n$ is constant. Then, for $r\ge r_1$, we~have $r'= (\ln r)^{2n}/4$. Hence, for $D(r)=\Vert 1\vee \Deg\Vert_{B(r)}$, we~obtain
\begin{align*}
n_o'(r)&=
n \Bigl(1+\frac{\ln(r')}{\ln(r/r')} +{280 n{\frac{\ln (r)}{\sqrt[ n]{r'}}}}\Bigr)\Bigl(1+ \frac{1}{2}\frac{\ln (D(r))} {\ln (r)}\Bigr)\\
&\le n \Bigl(1+\frac{2n \ln(\ln r) }{\ln(r / (\ln r)^{2n}) } +\frac{560 n}{\ln (r)}\Bigr)\Bigl(1+ \frac{1}{2}\frac{\ln (D(r))} {\ln (r)}\Bigr)\\
&\le n \Bigl(1+\frac{4n \ln(\ln r) }{\ln(r) } +\frac{560 n}{\ln (r)}\Bigr)\Bigl(1+ \frac{1}{2}\frac{\ln (D(r))} {\ln (r)}\Bigr)\\
&\le n \Bigl(1+\frac{284n \ln(\ln r) }{\ln(r) } \Bigr)\Bigl(1+ \frac{1}{2}\frac{\ln (D(r))} {\ln (r)}\Bigr),
\end{align*}
where we used that $(\ln r)^{4n} \le r$ implies $\ln(r/( \ln r)^{2n})\ge \ln(r)/2$ for $r\ge r_1$ in the second to last step and it also implies $2<\ln 8< \ln 4n <\ln \ln r$ which is used in the last step. This finishes the proof.
\end{proof}



\begin{thebibliography}{99}
\bibitem{AndresDS-16}S. Andres, J.-D. Deuschel \& M. Slowik - “Heat kernel estimates for random walks with degenerate weights”, Electron. J. Probab. 21 (2016), article ID 33, 21 pages | Zbl | MR | DOI
\bibitem{Barlow-04}M. T. Barlow - “Random walks on supercritical percolation clusters”, Ann. Probab. 32 (2004) no. 4, p. 3024-3084 | Zbl | MR | DOI
\bibitem{Barlow-book}M. T. Barlow - Random walks and heat kernels on graphs, London Math. Soc. Lect. Note Series, vol. 438, Cambridge University Press, Cambridge, 2017 | Zbl | MR | DOI
\bibitem{BarlowChen-16}M. T. Barlow \& X. Chen - “Gaussian bounds and parabolic Harnack inequality on locally irregular graphs”, Math. Ann. 366 (2016) no. 3-4, p. 1677-1720 | Zbl | MR | DOI
\bibitem{BakryCLS-95}D. Bakry, T. Coulhon, M. Ledoux \& L. Saloff-Coste - “Sobolev inequalities in disguise”, Indiana Univ. Math. J. 44 (1995) no. 4, p. 1033-1074 | Zbl | MR | DOI
\bibitem{BCS}S. Boutayeb, T. Coulhon \& A. Sikora - “A new approach to pointwise heat kernel upper bounds on doubling metric measure spaces”, Adv. Math. 270 (2015), p. 302-374 | Zbl | MR | DOI
\bibitem{BauerKH-13}F. Bauer, B. Hua \& M. Keller - “On the $\ell ^p$ spectrum of Laplacians on graphs”, Adv. Math. 248 (2013), p. 717-735 | Zbl | MR | DOI
\bibitem{BauerHuaYau-17}F. Bauer, B. Hua \& S.-T. Yau - “Sharp Davies-Gaffney-Grigor’yan lemma on graphs”, Math. Ann. 368 (2017) no. 3-4, p. 1429-1437 | Zbl | MR | DOI
\bibitem{BauerKW-15}F. Bauer, M. Keller \& R. K. Wojciechowski - “Cheeger inequalities for unbounded graph Laplacians”, J. Eur. Math. Soc. (JEMS) 17 (2015) no. 2, p. 259-271 | DOI | Zbl | MR
\bibitem{Bella-22}P. Bella \& M. Schäffner - “Non-uniformly parabolic equations and applications to the random conductance model”, Probab. Theory Relat. Fields 182 (2022) no. 1-2, p. 353-397 | Zbl | MR | DOI
\bibitem{CarlenKS-87}E. A. Carlen, S. Kusuoka \& D. W. Stroock - “Upper bounds for symmetric Markov transition functions”, Ann. Inst. H. Poincaré Probab. Statist. 23 (1987) no. 2, p. 245-287 | Numdam | Zbl | MR
\bibitem{ChenKW-20}X. Chen, T. Kumagai \& J. Wang - “Random conductance models with stable-like jumps: heat kernel estimates and Harnack inequalities”, J. Funct. Anal. 279 (2020) no. 7, p. 108656, 51 | Zbl | MR | DOI
\bibitem{Coulhon-92}T. Coulhon - “Inégalités de Gagliardo-Nirenberg pour les semi-groupes d’opérateurs et applications”, Potential Anal. 1 (1992) no. 4, p. 343-353 | Zbl | MR | DOI
\bibitem{Davies-93a}E. B. Davies - “Analysis on graphs and noncommutative geometry”, J. Funct. Anal. 111 (1993) no. 2, p. 398-430 | Zbl | MR | DOI
\bibitem{Davies-93}E. B. Davies - “Large deviations for heat kernels on graphs”, J. London Math. Soc. (2) 47 (1993) no. 1, p. 65-72 | Zbl | MR | DOI
\bibitem{Delmotte-97}T. Delmotte - “Inégalité de Harnack elliptique sur les graphes”, Colloq. Math. 72 (1997) no. 1, p. 19-37 | Zbl | MR | DOI
\bibitem{Delmotte-99}T. Delmotte - “Parabolic Harnack inequality and estimates of Markov chains on graphs”, Rev. Mat. Iberoamericana 15 (1999) no. 1, p. 181-232 | MR | Zbl | DOI
\bibitem{Folz-11}M. Folz - “Gaussian upper bounds for heat kernels of continuous time simple random walks”, Electron. J. Probab. 16 (2011), p. 1693-1722 | MR | Zbl | DOI
\bibitem{Folz-14b}M. Folz - “Volume growth and spectrum for general graph Laplacians”, Math. Z. 276 (2014) no. 1-2, p. 115-131 | MR | Zbl | DOI
\bibitem{Folz-14}M. Folz - “Volume growth and stochastic completeness of graphs”, Trans. Amer. Math. Soc. 366 (2014) no. 4, p. 2089-2119 | MR | Zbl | DOI
\bibitem{GrigoryanHH}A. Grigor’yan, E. Hu \& J. Hu - “Parabolic mean value inequality and on-diagonal upper bound of the heat kernel on doubling spaces”, Math. Ann. 389 (2024) no. 3, p. 2411-2467 | MR | Zbl | DOI
\bibitem{GrigoryanHuangMasamune-12}A. Grigor’yan, X. Huang \& J. Masamune - “On stochastic completeness of jump processes”, Math. Z. 271 (2012) no. 3-4, p. 1211-1239 | MR | Zbl | DOI
\bibitem{Grigoryan-94}A. Grigor’yan - “Heat kernel upper bounds on a complete non-compact manifold”, Rev. Mat. Iberoamericana 10 (1994) no. 2, p. 395-452 | MR | Zbl | DOI
\bibitem{HuangKS-20}X. Huang, M. Keller \& M. Schmidt - “On the uniqueness class, stochastic completeness and volume growth for graphs”, Trans. Amer. Math. Soc. 373 (2020) no. 12, p. 8861-8884 | MR | Zbl | DOI
\bibitem{HaeselerKW-13}S. Haeseler, M. Keller \& R. K. Wojciechowski - “Volume growth and bounds for the essential spectrum for Dirichlet forms”, J. London Math. Soc. (2) 88 (2013) no. 3, p. 883-898 | MR | Zbl | DOI
\bibitem{Keller-15}M. Keller - “Intrinsic metrics on graphs: a survey”, in Mathematical technology of networks, Springer Proc. Math. Stat., vol. 128, Springer, Cham, 2015, p. 81-119 | Zbl | DOI
\bibitem{KellerKNR-24}M. Keller, A. Kostenko, N. Nicolussi \& C. Rose - “Isoperimetry and heat kernels on path graphs”, In preparation
\bibitem{KellerLW-21}M. Keller, D. Lenz \& R. K. Wojciechowski - Graphs and discrete Dirichlet spaces, Grundlehren Math. Wissen., vol. 358, Springer, Cham, 2021 | MR | Zbl | DOI
\bibitem{KellerMuench}M. Keller \& F. Münch - “A new discrete Hopf-Rinow theorem”, Discrete Math. 342 (2019) no. 9, p. 2751-2757 | DOI | MR | Zbl
\bibitem{Kostenko-21}A. Kostenko - “Heat kernels of the discrete Laguerre operators”, Lett. Math. Phys. 111 (2021) no. 2, article ID 32, 29 pages | MR | DOI | Zbl
\bibitem{KellerRose-22b}M. Keller \& C. Rose - “Anchored heat kernel upper bounds on graphs with unbounded geometry and anti-trees”, Calc. Var. Partial Differential Equations 63 (2024) no. 1, article ID 20, 18 pages | DOI | MR | Zbl
\bibitem{KellerRose-22a}M. Keller \& C. Rose - “Gaussian upper bounds for heat kernels on graphs with unbounded geometry”, J. Spectral Theory 16 (2026) no. 2, p. 709-750 | DOI | MR | Zbl
\bibitem{Pang-93}M. M. H. Pang - “Heat kernels of graphs”, J. London Math. Soc. (2) 47 (1993) no. 1, p. 50-64 | DOI | MR | Zbl
\bibitem{SaloffCoste-92}L. Saloff-Coste - “A note on Poincaré, Sobolev, and Harnack inequalities”, Internat. Math. Res. Notices (1992) no. 2, p. 27-38 | DOI | MR | Zbl
\bibitem{SaloffCoste-92a}L. Saloff-Coste - “Uniformly elliptic operators on Riemannian manifolds”, J. Differential Geom. 36 (1992) no. 2, p. 417-450 | DOI | MR | Zbl
\bibitem{SaloffC-01}L. Saloff-Coste - Aspects of Sobolev-type inequalities, London Math.Soc. Lect. Note Series, vol. 289, Cambridge University Press, Cambridge, 2002 | MR | Zbl
\bibitem{Trudinger-71}N. S. Trudinger - “On the regularity of generalized solutions of linear, non-uniformly elliptic equations”, Arch. Rational Mech. Anal. 42 (1971), p. 50-62 | DOI | MR | Zbl
\bibitem{Varopoulos-85}N. T. Varopoulos - “Hardy-Littlewood theory for semigroups”, J. Funct. Anal. 63 (1985) no. 2, p. 240-260 | DOI
\end{thebibliography}
\end{document}
